\chapter{อัลกอริทึมควอนตัม คอมพิวติง และพรมแดนประยุกต์}
\label{chap:quantum_algorithms_frontiers}
\index{อัลกอริทึมควอนตัม (Quantum Algorithms)}
\index{ความได้เปรียบเชิงควอนตัม (Quantum Advantage)}
\index{อัลกอริทึมดอยช์-โยซา (Deutsch-Jozsa Algorithm)}
\index{อัลกอริทึมของโกรเวอร์ (Grover's Algorithm)}
\index{การแปลงฟูเรียร์ควอนตัม (Quantum Fourier Transform ---- QFT)}
\index{อัลกอริทึมของชอร์ (Shor's Algorithm)}
\index{การประเมินเฟสควอนตัม (Quantum Phase Estimation ---- QPE)}
\index{อัลกอริทึมแปรผันไอเกนควอนตัม (VQE)}
\index{รหัสแก้ความผิดพลาดควอนตัม (Quantum Error Correction)}
\index{คิวบิตตัวนำยิ่งยวด (Superconducting Qubits)}
\index{เซนเซอร์ควอนตัม (Quantum Sensing)}

% ==========================================================
% แผนบริหารการสอนประจำบทที่ 10
% ==========================================================
\section*{แผนบริหารการสอนประจำบทที่ 10}
\addcontentsline{toc}{section}{แผนบริหารการสอนประจำบทที่ 10}

\begin{lessonplanbox}
\noindent\textbf{หัวข้อเนื้อหาประจำบท}
\begin{enumerate}
    \item ความได้เปรียบเชิงควอนตัม (Quantum Advantage) และการประมวลผลขนานควอนตัม
    \item อัลกอริทึมของดอยช์และดอยช์-โยซา (Deutsch \& Deutsch-Jozsa Algorithms)
    \item อัลกอริทึมการค้นหาของโกรเวอร์เชิงเรขาคณิตและการขยายแอมพลิจูด
    \item การแปลงฟูเรียร์ควอนตัม (Quantum Fourier Transform ---- QFT)
    \item การประเมินเฟสควอนตัม (QPE) และอัลกอริทึมการแยกตัวประกอบของชอร์
    \item ขั้นตอนวิธีแบบผสมผสานคลาสสิก-ควอนตัมในยุค NISQ (VQE และ QAOA)
    \item ทฤษฎีรหัสแก้ความผิดพลาดควอนตัม (QEC) และเซอร์เฟสโค้ด (Surface Codes)
    \item สถาปัตยกรรมฮาร์ดแวร์ควอนตัมคอมพิวติง (ทรานสมอน ไอออนดักจับ และอะตอมเป็นกลาง)
    \item เซนเซอร์ควอนตัมและมาตรวิทยาควอนตัม (NV Center ในการเกษตรดิจิทัลและสิ่งแวดล้อม)
    \item ห้องปฏิบัติการเสมือนจริง 3 มิติ WebAR ชิปประมวลผลควอนตัมและโกรเวอร์
    \item การจำลองวงจรควอนตัมอัลกอริทึมด้วยภาษา Python และ Qiskit
    \item ตัวอย่างการแก้โจทย์เชิงวิเคราะห์ตามระเบียบวิธี P-S-C-T
    \item สรุปสาระสำคัญและแบบฝึกหัดท้ายบท
\end{enumerate}

\noindent\textbf{วัตถุประสงค์เชิงพฤติกรรม}
\begin{enumerate}
    \item อธิบายกลไกการแทรกสอดเชิงสร้างสรรค์และทำลายล้างที่สร้างความได้เปรียบเชิงควอนตัมเหนือคอมพิวเตอร์คลาสสิกได้อย่างลึกซึ้ง
    \item คำนวณและแสดงการวิวัฒน์ของเวกเตอร์สถานะในอัลกอริทึมของดอยช์และดอยช์-โยซาผ่านกลไกการสะท้อนเฟสได้อย่างแม่นยำ
    \item วิเคราะห์การทำงานของออราเคิลและตัวดำเนินการแพร่กระจายของโกรเวอร์ในระนาบสองมิติและคำนวณจำนวนรอบทำซ้ำที่เหมาะสมได้
    \item เขียนและอธิบายวงจรการแปลงฟูเรียร์ควอนตัม (QFT) พร้อมคำนวณการแปลงสถานะฐานได้อย่างถูกต้อง
    \item อธิบายขั้นตอนการลดรูปปัญหาการแยกตัวประกอบเฉพาะสู่การหาคาบในอัลกอริทึมของชอร์และวิเคราะห์ผลกระทบต่อวิทยาการรหัสลับได้
    \item อธิบายหลักการทำงานของอัลกอริทึมแบบผสมผสาน VQE ในการหาพลังงานสถานะพื้นของโมเลกุลในยุคฮาร์ดแวร์ NISQ ได้
    \item วิเคราะห์กลไกการตรวจจับความผิดพลาดบิตฟลิปและเฟสฟลิปในรหัสแก้ความผิดพลาดควอนตัมและหลักการทำงานของเซอร์เฟสโค้ดได้
    \item เปรียบเทียบข้อดี ข้อจำกัด และตัวแปรทางฟิสิกส์ของสถาปัตยกรรมฮาร์ดแวร์ควอนตัมชั้นนำของโลกได้อย่างมีวิจารณญาณ
    \item ประยุกต์ใช้เซนเซอร์ควอนตัมจุดบกพร่องไนโตรเจน-ช่องว่าง (NV Center) ในการตรวจวัดสนามแม่เหล็กและธาตุอาหารในดินทางการเกษตรได้
    \item พัฒนาโปรแกรมภาษา Python ร่วมกับไลบรารี Qiskit ในการจำลองอัลกอริทึมของโกรเวอร์และประเมินผลความน่าจะเป็นได้อย่างสมบูรณ์
\end{enumerate}
\end{lessonplanbox}

% ==========================================================
% 10.1 ความได้เปรียบเชิงควอนตัมและการขนานเชิงควอนตัม
% ==========================================================
\newpage
\section{ความได้เปรียบเชิงควอนตัมและการประมวลผลขนาน}
\index{ความได้เปรียบเชิงควอนตัม (Quantum Advantage)}
\index{การขนานเชิงควอนตัม (Quantum Parallelism)}
\index{คลาสความซับซ้อน BQP}
\index{การแทรกสอดเชิงควอนตัม (Quantum Interference)}

เป้าหมายสูงสุดของการพัฒนาเทคโนโลยีสารสนเทศควอนตัมคือการบรรลุ \textbf{ความได้เปรียบเชิงควอนตัม (Quantum Advantage)} หรือที่เคยได้รับการขนานนามในอดีตว่า อภิมหาอำนาจควอนตัม (Quantum Supremacy) ซึ่งหมายถึงสภาวะที่คอมพิวเตอร์ควอนตัมสามารถประมวลผลเพื่อแก้ไขปัญหาคำนวณที่มีความซับซ้อนสูงได้เร็วกว่าซูเปอร์คอมพิวเตอร์คลาสสิกที่ทรงพลังที่สุดในโลกอย่างมีนัยสำคัญยิ่งยวด โดยปัญหาดังกล่าวอาจต้องใช้เวลาประมวลผลบนซูเปอร์คอมพิวเตอร์แบบดั้งเดิมเป็นเวลาหลายพันปี แต่สามารถแก้ไขได้บนระบบควอนตัมภายในเวลาเพียงไม่กี่นาที

\subsection{คลาสความซับซ้อนเชิงคำนวณและบทบาทของ BQP}
ในทฤษฎีความซับซ้อนของการคำนวณ (Computational Complexity Theory) ปัญหาเชิงการตัดสินใจถูกจัดกลุ่มเป็นคลาสต่างๆ ตามทรัพยากรเวลาและหน่วยความจำที่ต้องใช้
\begin{itemize}
    \item \textbf{คลาส P (Polynomial Time)} ---- ปัญหาที่สามารถแก้ไขได้ด้วยเครื่องจักรทัวริงดั้งเดิมเชิงกำหนด (Deterministic Turing Machine) ภายในเวลาพหุนาม $O(n^k)$
    \item \textbf{คลาส NP (Nondeterministic Polynomial Time)} ---- ปัญหาที่เมื่อได้รับคำตอบที่เสนอมาแล้ว สามารถตรวจสอบความถูกต้องได้ภายในเวลาพหุนาม
    \item \textbf{คลาส BPP (Bounded-error Probabilistic Polynomial Time)} ---- ปัญหาที่สามารถแก้ไขได้ด้วยคอมพิวเตอร์คลาสสิกที่ใช้อัลกอริทึมเชิงสุ่ม โดยมีความน่าจะเป็นของความผิดพลาดไม่เกิน $1/3$
    \item \textbf{คลาส BQP (Bounded-error Quantum Polynomial Time)} ---- ปัญหาการตัดสินใจที่สามารถแก้ไขได้ด้วยคอมพิวเตอร์ควอนตัมภายในเวลาพหุนาม โดยมีความน่าจะเป็นของความผิดพลาดไม่เกิน $1/3$
\end{itemize}

เป็นที่ทราบกันดีว่า $\text{P} \subseteq \text{BPP} \subseteq \text{BQP}$ ทว่าความสัมพันธ์ระหว่าง $\text{BQP}$ และ $\text{NP}$ ยังคงเป็นหัวข้อการวิจัยที่สำคัญยิ่งในปัจจุบัน นักวิทยาศาสตร์เชื่อว่า $\text{BQP}$ ครอบคลุมปัญหาบางประเภทที่อยู่นอกเหนือคลาส $\text{BPP}$ แต่อาจไม่ครอบคลุมปัญหา $\text{NP-Complete}$ ทั้งหมด ดังแสดงในความสามารถในการแยกตัวประกอบจำนวนเต็มของอัลกอริทึมชอร์ ซึ่งอยู่ใน $\text{BQP}$ แต่ไม่เชื่อว่าอยู่ใน $\text{P}$ หรือ $\text{BPP}$

\subsection{การประมวลผลขนานเชิงควอนตัมและการแทรกสอด}
ขุมพลังที่ขับเคลื่อนความเร็วอันมหาศาลของคอมพิวเตอร์ควอนตัมประกอบด้วยเสาหลักสองประการที่ต้องทำงานประสานกันอย่างลงตัว

\noindent\textbf{1. การขนานเชิงควอนตัม (Quantum Parallelism)} \\
พิจารณาระบบ $n$ คิวบิตที่เริ่มต้นในสถานะฐานศูนย์ทั้งหมด $\ket{0}^{\otimes n}$ เมื่อส่งคิวบิตทั้งหมดผ่านเกตฮาดามาร์ด $H^{\otimes n}$ สถานะของระบบจะกลายเป็นการซ้อนทับเชิงเส้นที่มีขนาดแอมพลิจูดสม่ำเสมอของสถานะฐานที่เป็นไปได้ทั้งหมด $2^n$ สถานะ
\begin{equation}
\ket{\psi_0} = H^{\otimes n} \ket{0}^{\otimes n} = \frac{1}{\sqrt{2^n}} \sum_{x=0}^{2^n-1} \ket{x}
\label{eq:hadamard_uniform_superposition}
\end{equation}
หากเรามีฟังก์ชันทางคณิตศาสตร์ $f: \{0, 1\}^n \to \{0, 1\}^m$ ซึ่งสร้างขึ้นในรูปของตัวดำเนินการเอกภาพ $U_f$ ที่กระทำบนสองรีจิสเตอร์ตามสมการ
\begin{equation}
U_f \big( \ket{x} \otimes \ket{0}^{\otimes m} \big) = \ket{x} \otimes \ket{f(x)}
\end{equation}
เมื่อเราป้อนสถานะซ้อนทับ $\ket{\psi_0}$ เข้าสู่วงจร จะได้ผลลัพธ์เป็น
\begin{equation}
U_f \left( \frac{1}{\sqrt{2^n}} \sum_{x=0}^{2^n-1} \ket{x} \otimes \ket{0}^{\otimes m} \right) = \frac{1}{\sqrt{2^n}} \sum_{x=0}^{2^n-1} \ket{x} \otimes \ket{f(x)}
\label{eq:quantum_parallel_evaluation}
\end{equation}
สมการ (\ref{eq:quantum_parallel_evaluation}) แสดงให้เห็นว่า การประเมินค่าฟังก์ชัน $f(x)$ สำหรับค่าอินพุตที่เป็นไปได้ทั้งหมด $2^n$ ค่า ถูกคำนวณขึ้นพร้อมกันในการกระทำของเกตเพียงครั้งเดียว! สำหรับ $n = 50$ คิวบิต จำนวนสถานะที่ประมวลผลพร้อมกันจะมีค่าสูงถึง $2^{50} \approx 1.125 \times 10^{15}$ สถานะ

\noindent\textbf{2. การแทรกสอดเชิงควอนตัม (Quantum Interference)} \\
แม้ว่าการขนานเชิงควอนตัมจะประเมินค่าฟังก์ชันได้พร้อมกันทุกค่า ทว่าตามกฎของบอร์น หากเราทำการวัดรีจิสเตอร์โดยตรง ระบบจะยุบตัวลงสู่สถานะคู่ลำดับ $\ket{x_0}\ket{f(x_0)}$ สุ่มเพียงค่าเดียวเท่านั้น ซึ่งไม่ได้ให้ความได้เปรียบเหนือการสุ่มตรวจในคอมพิวเตอร์คลาสสิกเลย

หัวใจอันแยบยลของอัลกอริทึมควอนตัมจึงอยู่ที่การใช้ \textbf{การแทรกสอดเชิงควอนตัม} ผ่านลำดับของเกตเอกภาพก่อนการวัด เพื่อทำให้แอมพลิจูดความน่าจะเป็นของสถานะคำตอบที่ไม่ต้องการเกิด \textbf{การแทรกสอดแบบหักล้าง (Destructive Interference)} จนแอมพลิจูดลดลงเหลือศูนย์ ในขณะที่แอมพลิจูดของสถานะคำตอบที่ถูกต้องเกิด \textbf{การแทรกสอดแบบเสริมสร้าง (Constructive Interference)} จนพุ่งสูงขึ้นเข้าใกล้หนึ่ง ทำให้เมื่อทำการวัดในขั้นตอนสุดท้าย ผู้ใช้งานจะได้คำตอบที่ถูกต้องด้วยความน่าจะเป็นที่สูงยิ่ง

% ==========================================================
% 10.2 อัลกอริทึมดอยช์และดอยช์-โยซา
% ==========================================================
\section{อัลกอริทึมของดอยช์และดอยช์-โยซา}
\index{อัลกอริทึมของดอยช์ (Deutsch Algorithm)}
\index{อัลกอริทึมดอยช์-โยซา (Deutsch-Jozsa Algorithm)}
\index{การสะท้อนเฟส (Phase Kickback)}
\index{ฟังก์ชันคงตัวและฟังก์ชันสมดุล}

ในคริสต์ศักราช 1985 เดวิด ดอยช์ (David Deutsch) แห่งมหาวิทยาลัยออกซฟอร์ด ได้เสนออัลกอริทึมควอนตัมแบบแรกของโลกที่พิสูจน์ว่าคอมพิวเตอร์ควอนตัมสามารถทำงานได้เร็วกว่าคอมพิวเตอร์คลาสสิก ต่อมาในปี ค.ศ. 1992 ดอยช์ ร่วมกับ ริชาร์ด โยซา (Richard Jozsa) ได้ขยายผลสู่อัลกอริทึมสำหรับ $n$ คิวบิต ซึ่งแสดงความได้เปรียบแบบก้าวกระโดดเชิงกำหนด (Deterministic Exponential Speedup)

\subsection{นิยามของปัญหาดอยช์-โยซา}
กำหนดให้ออราเคิลคำนวณฟังก์ชันบูลีน $f: \{0, 1\}^n \to \{0, 1\}$ โดยมีเงื่อนไขรับประกันว่าฟังก์ชัน $f$ จะต้องเป็นฟังก์ชันประเภทใดประเภทหนึ่งในสองประเภทนี้เท่านั้น
\begin{enumerate}
    \item \textbf{ฟังก์ชันคงตัว (Constant Function)} ---- ให้ค่าผลลัพธ์ $f(x)$ เป็นค่าเดียวกันทั้งหมด ไม่ว่าอินพุต $x$ จะเป็นอะไร (นั่นคือ เป็น $0$ ทั้งหมด หรือเป็น $1$ ทั้งหมด)
    \item \textbf{ฟังก์ชันสมดุล (Balanced Function)} ---- ให้ค่าผลลัพธ์ $f(x) = 0$ สำหรับอินพุตจำนวนครึ่งหนึ่ง ($2^{n-1}$ ค่า) และให้ค่า $f(x) = 1$ สำหรับอินพุตอีกครึ่งหนึ่งที่เหลือ ($2^{n-1}$ ค่า)
\end{enumerate}
\textbf{เป้าหมาย} ---- จงระบุให้ได้อย่างแน่นอนเด็ดขาดว่า ฟังก์ชัน $f$ เป็นฟังก์ชันคงตัวหรือฟังก์ชันสมดุล โดยการเรียกใช้งานออราเคิลให้น้อยครั้งที่สุด

ในคอมพิวเตอร์คลาสสิก ในกรณีแย่ที่สุด (Worst-case Scenario) หากเราสุ่มตรวจค่าอินพุตไปแล้ว $2^{n-1}$ ค่าและพบว่าเป็น $0$ เหมือนกันทั้งหมด เรายังไม่สามารถสรุปได้ว่าฟังก์ชันเป็นแบบคงตัว เพราะค่าที่ $2^{n-1}+1$ อาจเป็น $1$ ก็ได้ ดังนั้น อัลกอริทึมคลาสสิกแบบกำหนดผลแน่นอนจะต้องเรียกใช้ออราเคิลมากถึง
\begin{equation}
N_{\text{classical}} = 2^{n-1} + 1 \text{ ครั้ง}
\end{equation}
สำหรับ $n = 32$ จะต้องเรียกใช้มากกว่าสองพันล้านครั้ง!

ทว่าอัลกอริทึมดอยช์-โยซา สามารถหาคำตอบได้อย่างถูกต้องแน่นอนร้อยเปอร์เซ็นต์ โดยการเรียกใช้ออราเคิล \textbf{เพียง 1 ครั้งเท่านั้น ($N_{\text{quantum}} = 1$)}

\subsection{กลไกการสะท้อนเฟส (Phase Kickback)}
กุญแจสำคัญของอัลกอริทึมคือการเตรียมคิวบิตช่วย (Ancilla Qubit) ในสถานะ $\ket{-} = \frac{\ket{0}-\ket{1}}{\sqrt{2}}$ พิจารณาการกระทำของเกตออราเคิล $U_f$ ต่อสถานะ $\ket{x}\ket{-}$
\begin{align}
U_f \big( \ket{x}\ket{-} \big) &= U_f \left( \ket{x} \otimes \frac{\ket{0}-\ket{1}}{\sqrt{2}} \right) \nonumber \\
&= \frac{\ket{x}\ket{0 \oplus f(x)} - \ket{x}\ket{1 \oplus f(x)}}{\sqrt{2}}
\end{align}
\begin{itemize}
    \item หาก $f(x) = 0$ จะได้ $\frac{\ket{x}\ket{0} - \ket{x}\ket{1}}{\sqrt{2}} = (+1) \ket{x}\ket{-}$
    \item หาก $f(x) = 1$ จะได้ $\frac{\ket{x}\ket{1} - \ket{x}\ket{0}}{\sqrt{2}} = (-1) \ket{x}\ket{-}$
\end{itemize}
ดังนั้น ในทุกกรณี เราสามารถเขียนผลลัพธ์ได้อย่างกระชับเป็น
\begin{equation}
U_f \big( \ket{x}\ket{-} \big) = (-1)^{f(x)} \ket{x}\ket{-}
\label{eq:phase_kickback_identity}
\end{equation}
ข้อมูลของฟังก์ชัน $f(x)$ ถูกถ่ายโอน (Kickback) ขึ้นมาเป็นเฟสสัมพัทธ์ของคิวบิตอินพุต $\ket{x}$ โดยที่คิวบิตช่วยยังคงอยู่ในสถานะ $\ket{-}$ เช่นเดิมไม่เปลี่ยนแปลง

\subsection{การวิวัฒน์ของสถานะในอัลกอริทึมดอยช์-โยซา}
ขั้นตอนการทำงานของอัลกอริทึมดอยช์-โยซาประกอบด้วยลำดับการแปลงสถานะดังนี้
\begin{enumerate}
    \item \textbf{การเตรียมสถานะเริ่มต้น} ---- เตรียม $n$ คิวบิตแรกในสถานะ $\ket{0}$ และคิวบิตช่วยในสถานะ $\ket{1}$
    \begin{equation}
    \ket{\psi_0} = \ket{0}^{\otimes n} \otimes \ket{1}
    \end{equation}
    \item \textbf{การสร้างสถานะซ้อนทับสม่ำเสมอ} ---- นำเกตฮาดามาร์ดกระทำกับทุกคิวบิต
    \begin{equation}
    \ket{\psi_1} = (H^{\otimes n} \otimes H)\ket{\psi_0} = \left( \frac{1}{\sqrt{2^n}} \sum_{x=0}^{2^n-1} \ket{x} \right) \otimes \ket{-}
    \end{equation}
    \item \textbf{การประเมินออราเคิลและการสะท้อนเฟส} ---- ส่งสถานะผ่านออราเคิล $U_f$ โดยใช้สมบัติ (\ref{eq:phase_kickback_identity})
    \begin{equation}
    \ket{\psi_2} = U_f \ket{\psi_1} = \left( \frac{1}{\sqrt{2^n}} \sum_{x=0}^{2^n-1} (-1)^{f(x)} \ket{x} \right) \otimes \ket{-}
    \end{equation}
    \item \textbf{การแทรกสอดด้วยการแปลงฮาดามาร์ดครั้งสุดท้าย} ---- กระทำเกต $H^{\otimes n}$ บน $n$ คิวบิตแรก โดยใช้เอกลักษณ์ $H^{\otimes n}\ket{x} = \frac{1}{\sqrt{2^n}}\sum_{y=0}^{2^n-1} (-1)^{x \cdot y} \ket{y}$ โดยที่ $x \cdot y = x_1 y_1 \oplus \dots \oplus x_n y_n$
    \begin{align}
    \ket{\psi_3} &= (H^{\otimes n} \otimes I) \ket{\psi_2} \nonumber \\
    &= \frac{1}{2^n} \sum_{y=0}^{2^n-1} \left( \sum_{x=0}^{2^n-1} (-1)^{f(x) \oplus (x \cdot y)} \right) \ket{y} \otimes \ket{-}
    \label{eq:deutsch_jozsa_final_expansion}
    \end{align}
\end{enumerate}

พิจารณาแอมพลิจูดความน่าจะเป็นของสถานะฐานศูนย์ทั้งหมด $\ket{0}^{\otimes n}$ (นั่นคือ กรณี $y = 0$ ซึ่งทำให้ $x \cdot y = 0$ สำหรับทุก $x$)
\begin{equation}
C_{0^{\otimes n}} = \frac{1}{2^n} \sum_{x=0}^{2^n-1} (-1)^{f(x)}
\end{equation}
\begin{itemize}
    \item \textbf{กรณีฟังก์ชันคงตัว} ---- หาก $f(x) = c \in \{0, 1\}$ สำหรับทุก $x$
    \begin{equation}
    C_{0^{\otimes n}} = \frac{1}{2^n} \sum_{x=0}^{2^n-1} (-1)^c = \frac{1}{2^n} \big( 2^n (-1)^c \big) = \pm 1
    \end{equation}
    ความน่าจะเป็นในการวัดได้สถานะ $\ket{0}^{\otimes n}$ คือ $P(0^{\otimes n}) = |C_{0^{\otimes n}}|^2 = 1$ (ร้อยละ 100 เต็ม)
    \item \textbf{กรณีฟังก์ชันสมดุล} ---- เนื่องจากมี $x$ ที่ทำให้ $f(x)=0$ อยู่ $2^{n-1}$ ค่า และมี $x$ ที่ทำให้ $f(x)=1$ อยู่ $2^{n-1}$ ค่า ผลรวมเครื่องหมายจะหักล้างกันหมดสิ้น
    \begin{equation}
    C_{0^{\otimes n}} = \frac{1}{2^n} \Big( 2^{n-1}(+1) + 2^{n-1}(-1) \Big) = 0
    \end{equation}
    ความน่าจะเป็นในการวัดได้สถานะ $\ket{0}^{\otimes n}$ คือ $P(0^{\otimes n}) = 0$
\end{itemize}

ดังนั้น เมื่อทำการวัด $n$ คิวบิตแรก หากผลลัพธ์เป็น $\ket{0}^{\otimes n}$ เราสามารถสรุปได้ทันทีว่าฟังก์ชันเป็น \textbf{ฟังก์ชันคงตัว} แต่หากได้ผลลัพธ์เป็นสถานะอื่นใดที่ไม่ใช่ศูนย์ทั้งหมด สรุปได้ทันทีว่าเป็น \textbf{ฟังก์ชันสมดุล} ด้วยการเรียกใช้ออราเคิลเพียงครั้งเดียวเท่านั้น!

% ==========================================================
% 10.3 อัลกอริทึมการค้นหาของโกรเวอร์
% ==========================================================
\section{อัลกอริทึมการค้นหาของโกรเวอร์เชิงเรขาคณิต}
\index{อัลกอริทึมของโกรเวอร์ (Grover's Algorithm)}
\index{การขยายแอมพลิจูด (Amplitude Amplification)}
\index{การหมุนในระนาบสองมิติ}
\index{อัตราเร่งแบบกำลังสอง (Quadratic Speedup)}

ในคริสต์ศักราช 1996 เลิฟ โกรเวอร์ (Lov Grover) นักวิทยาศาสตร์แห่งเบลล์แล็บส์ ได้ค้นพบอัลกอริทึมสำหรับการค้นหาข้อมูลในฐานข้อมูลที่ไม่มีการจัดหมวดหมู่ (Unsorted Database Search) ซึ่งเป็นปัญหาพื้นฐานที่สุดปัญหาหนึ่งในวิทยาการคอมพิวเตอร์

สำหรับฐานข้อมูลขนาด $N = 2^n$ รายการ ที่มีข้อมูลเป้าหมายที่ต้องการค้นหาอยู่ $M$ รายการ ($1 \le M \ll N$) คอมพิวเตอร์ดั้งเดิมต้องทำการตรวจสอบทีละรายการ ซึ่งในกรณีเฉลี่ยต้องใช้เวลาในการค้นหา $O(N/M)$ การดำเนินการ ทว่าอัลกอริทึมของโกรเวอร์สามารถค้นหาข้อมูลเป้าหมายได้ภายในเวลาเพียง
\begin{equation}
R \approx \frac{\pi}{4}\sqrt{\frac{N}{M}} = O\left(\sqrt{\frac{N}{M}}\right)
\label{eq:grover_speedup_scaling}
\end{equation}
ซึ่งให้อัตราเร่งเชิงความเร็วในระดับยกกำลังสอง (Quadratic Speedup)

\subsection{การวิเคราะห์เชิงเรขาคณิตในปริภูมิย่อยสองมิติ}
ความงดงามของอัลกอริทึมโกรเวอร์อยู่ที่การที่การดำเนินการทั้งหมดสามารถจำลองได้อย่างสมบูรณ์แบบในรูปของการหมุนเวกเตอร์ในระนาบสองมิติจริง นิยามสถานะฐานตั้งฉากปรกติสองสถานะในปริภูมิฮิลเบิร์ต ได้แก่
\begin{align}
\ket{\beta} &= \frac{1}{\sqrt{M}} \sum_{x \in W} \ket{x} \qquad \text{(สถานะซ้อนทับสม่ำเสมอของคำตอบที่ถูกต้อง)} \label{eq:grover_target_state} \\
\ket{\alpha} &= \frac{1}{\sqrt{N-M}} \sum_{x \notin W} \ket{x} \qquad \text{(สถานะซ้อนทับสม่ำเสมอของคำตอบที่ไม่ถูกต้อง)} \label{eq:grover_nontarget_state}
\end{align}
โดยที่ $W$ คือเซตของดัชนีคำตอบเป้าหมาย สังเกตว่า $\braket{\alpha|\beta} = 0$ สถานะทั้งสองจึงก่อตัวเป็นฐานตั้งฉากของระนาบสองมิติ

สถานะเริ่มต้นหลังผ่านเกตฮาดามาร์ด $\ket{s} = \frac{1}{\sqrt{N}}\sum_{x=0}^{N-1}\ket{x}$ สามารถเขียนกระจายบนฐาน $\{\ket{\alpha}, \ket{\beta}\}$ ได้เป็น
\begin{equation}
\ket{s} = \sqrt{\frac{N-M}{N}} \ket{\alpha} + \sqrt{\frac{M}{N}} \ket{\beta} = \cos\left(\frac{\theta}{2}\right)\ket{\alpha} + \sin\left(\frac{\theta}{2}\right)\ket{\beta}
\label{eq:grover_initial_state_angle}
\end{equation}
โดยที่มุม $\theta$ นิยามจาก
\begin{equation}
\sin\left(\frac{\theta}{2}\right) = \sqrt{\frac{M}{N}} \implies \theta \approx 2\sqrt{\frac{M}{N}} \quad (\text{เมื่อ } M \ll N)
\end{equation}

\subsection{ตัวดำเนินการทำซ้ำของโกรเวอร์ (Grover Iteration)}
ในแต่ละรอบของการทำซ้ำ อัลกอริทึมจะประยุกต์ใช้ตัวดำเนินการรวม $G$ ซึ่งประกอบด้วยสองขั้นตอนย่อย
\begin{equation}
G = U_s U_w
\label{eq:grover_operator_def}
\end{equation}
\begin{enumerate}
    \item \textbf{ออราเคิลกลับเฟสเป้าหมาย (Phase Oracle ---- $U_w$)}
    \begin{equation}
    U_w = I - 2\sum_{x \in W} \ket{x}\bra{x} = I - 2\ket{\beta}\bra{\beta}
    \end{equation}
    ทางเรขาคณิต $U_w$ ทำหน้าที่สะท้อนเวกเตอร์สถานะข้ามแกน $\ket{\alpha}$ (กลับทิศทางองค์ประกอบตามแนวแกน $\ket{\beta}$)
    \item \textbf{ตัวดำเนินการแพร่กระจาย (Diffusion Operator ---- $U_s$)}
    \begin{equation}
    U_s = 2\ket{s}\bra{s} - I
    \end{equation}
    ทางเรขาคณิต $U_s$ ทำหน้าที่สะท้อนเวกเตอร์สถานะข้ามแนวเวกเตอร์ $\ket{s}$ ซึ่งสมนัยกับ \textbf{การพลิกแอมพลิจูดรอบค่าเฉลี่ย (Inversion about the Average)}
\end{enumerate}

\begin{figure}[H]
\centering
\begin{tikzpicture}[scale=1.1]
    % Wires
    \draw[line width=1.4pt, darkNavy] (-0.5, 1.2) -- (7.5, 1.2) node[right, font=\bfseries\color{darkNavy}] {$q_0$};
    \draw[line width=1.4pt, darkNavy] (-0.5, 0.0) -- (7.5, 0.0) node[right, font=\bfseries\color{darkNavy}] {$q_1$};
    
    % Initial States
    \node at (-1.0, 1.2) [font=\large\bfseries\color{cyberBlue}] {$\ket{0}$};
    \node at (-1.0, 0.0) [font=\large\bfseries\color{cyberBlue}] {$\ket{0}$};
    
    % First H Gates
    \draw[fill=rbruLightBg, draw=cyberBlue, line width=1.5pt, rounded corners=3pt] (0.5, 0.7) rectangle (1.5, 1.7);
    \node at (1.0, 1.2) [font=\large\bfseries\color{darkNavy}] {$H$};
    \draw[fill=rbruLightBg, draw=cyberBlue, line width=1.5pt, rounded corners=3pt] (0.5, -0.5) rectangle (1.5, 0.5);
    \node at (1.0, 0.0) [font=\large\bfseries\color{darkNavy}] {$H$};
    
    % Oracle Box
    \draw[fill=neonPurple!20!white, draw=neonPurple, line width=1.6pt, rounded corners=4pt] (2.2, -0.6) rectangle (4.0, 1.8);
    \node at (3.1, 0.6) [align=center, font=\small\bfseries\color{darkNavy}] {ออราเคิล\\$U_w$};
    
    % Diffuser Box
    \draw[fill=emeraldGlow!20!white, draw=emeraldGlow, line width=1.6pt, rounded corners=4pt] (4.6, -0.6) rectangle (6.8, 1.8);
    \node at (5.7, 0.6) [align=center, font=\small\bfseries\color{darkNavy}] {ดิฟฟิวเซอร์\\$2\ket{s}\bra{s}-I$};
\end{tikzpicture}
\caption{ผังวงจรควอนตัมอัลกอริทึมของโกรเวอร์ (เป้าหมายติดตามภาพ AR Image Target)}
\label{fig:grover_circuit_ar_target}
\end{figure}

% ==========================================================
% กล่องปฏิบัติการ WebAR & Image Tracking ประจำบทที่ 10
% ==========================================================
\begin{qrcodelabbox}{การทดลองอัลกอริทึมควอนตัมและสถาปัตยกรรมชิป WebAR}{https://tsanaphysics.github.io/ar-Quantum/lab10_grover_hardware.html}
\begin{minipage}[c]{0.70\textwidth}
\textbf{คำแนะนำการทดลองเสมือนจริงสามมิติ (Quantum Circuit Simulator)}
\begin{enumerate}
    \item สแกนรหัสตอบสนองเร็ว (QR Code) ด้านขวา เพื่อเปิดโปรแกรมจำลองวงจรโกรเวอร์ 3 มิติ
    \item ส่องกล้องเข้าหา \textbf{ภาพที่ \ref{fig:grover_circuit_ar_target}} ตัวจำลองจะเรนเดอร์โครงสร้างคิวบิตตัวนำยิ่งยวดแบบ Transmon Qubit พร้อมแสดงการขยายแอมพลิจูด (Amplitude Amplification)
    \item สังเกตการเพิ่มขึ้นของเวกเตอร์ความน่าจะเป็นของสถานะเป้าหมายในแต่ละรอบการทำซ้ำของโกรเวอร์
\end{enumerate}
\end{minipage}
\hfill
\begin{minipage}[c]{0.25\textwidth}
\centering
\qrcode[height=2.8cm]{https://tsanaphysics.github.io/ar-Quantum/lab10_grover_hardware.html}
\vspace{4pt}
\scriptsize\textbf{สแกนเพื่อเปิด WebAR}
\end{minipage}
\end{qrcodelabbox}

การสะท้อนสองครั้งติดต่อกันข้ามเส้นตรงสองเส้นที่ทำมุมกัน $\theta/2$ ย่อมเทียบเท่ากับ \textbf{การหมุนเวกเตอร์ในระนาบด้วยมุม $\theta$ ทวนเข็มนาฬิกา} มุ่งตรงไปยังแกนเป้าหมาย $\ket{\beta}$ หลังจากทำซ้ำ $k$ รอบ เวกเตอร์สถานะจะกลายเป็น
\begin{equation}
G^k \ket{s} = \cos\left( \frac{2k+1}{2}\theta \right) \ket{\alpha} + \sin\left( \frac{2k+1}{2}\theta \right) \ket{\beta}
\label{eq:grover_state_k_rotations}
\end{equation}

เพื่อให้ความน่าจะเป็นในการวัดได้สถานะเป้าหมาย $\ket{\beta}$ มีค่าสูงสุดเข้าใกล้ 1 เราต้องเลือกจำนวนรอบการทำซ้ำ $R$ ที่ทำให้มุมรวมเข้าใกล้ $\pi/2$ มากที่สุด
\begin{equation}
\frac{2R+1}{2}\theta \approx \frac{\pi}{2} \implies 2R \theta \approx \pi \implies R \approx \frac{\pi}{2\theta} \approx \frac{\pi}{4}\sqrt{\frac{N}{M}}
\end{equation}
หากผู้ใช้งานทำซ้ำเกินกว่า $R$ รอบ เวกเตอร์จะหมุนเลยแกน $\ket{\beta}$ ส่งผลให้ความน่าจะเป็นลดลง ซึ่งปรากฏการณ์นี้เรียกว่า การหมุนเกิน (Over-rotation)

% ==========================================================
% 10.4 การแปลงฟูเรียร์ควอนตัม
% ==========================================================
\section{การแปลงฟูเรียร์ควอนตัม (Quantum Fourier Transform ---- QFT)}
\index{การแปลงฟูเรียร์ควอนตัม (Quantum Fourier Transform)}
\index{การแปลงฟูเรียร์แบบเร็ว (FFT)}
\index{เกตเปลี่ยนเฟสควบคุม}

การแปลงฟูเรียร์แบบไม่ต่อเนื่อง (Discrete Fourier Transform ---- DFT) เป็นหนึ่งในเครื่องมือทางคณิตศาสตร์ที่มีการใช้งานกว้างขวางที่สุดในวิทยาศาสตร์และวิศวกรรม ในระบบควอนตัม การแปลงฟูเรียร์แบบไม่ต่อเนื่องถูกนำมาสร้างเป็นตัวดำเนินการเอกภาพบนปริภูมิฮิลเบิร์ต เรียกว่า \textbf{การแปลงฟูเรียร์ควอนตัม (Quantum Fourier Transform หรือ QFT)}

สำหรับระบบ $n$ คิวบิต ($N = 2^n$) QFT ทำหน้าที่แปลงสถานะฐานเชิงคำนวณ $\ket{j}$ ไปสู่การซ้อนทับเชิงเส้นของสถานะฐานทั้งหมดตามสมการ
\begin{equation}
\text{QFT} \ket{j} = \frac{1}{\sqrt{N}} \sum_{k=0}^{N-1} \omega^{j k} \ket{k} = \frac{1}{\sqrt{2^n}} \sum_{k=0}^{2^n-1} e^{2\pi i j k / 2^n} \ket{k}
\label{eq:qft_definition_sum}
\end{equation}
โดยที่ $\omega = e^{2\pi i / N}$ คือรากที่ $N$ ของหนึ่ง

\subsection{รูปแบบผลคูณเทนเซอร์ของ QFT}
ความก้าวล้ำที่ทำให้ QFT สามารถสร้างขึ้นได้ด้วยวงจรควอนตัมที่มีประสิทธิภาพสูง ถูกค้นพบโดย ดอน คอปเปอร์สมิธ (Don Coppersmith) ในปี ค.ศ. 1994 โดยแสดงว่าสถานะผลลัพธ์ในสมการ (\ref{eq:qft_definition_sum}) สามารถแยกตัวประกอบให้อยู่ในรูปผลคูณเทนเซอร์ของแต่ละคิวบิตได้อย่างสมบูรณ์

เขียนเลขจำนวนเต็ม $j$ ในรูปเลขฐานสอง $j = j_1 j_2 \dots j_n = j_1 2^{n-1} + j_2 2^{n-2} + \dots + j_n 2^0$ และแทนเศษส่วนฐานสองด้วยสัญลักษณ์ $0.j_l j_{l+1} \dots j_n = \sum_{m=l}^n j_m 2^{-(m-l+1)}$ จะได้
\begin{align}
\text{QFT} \ket{j_1 j_2 \dots j_n} = \frac{1}{\sqrt{2^n}} &\Big( \ket{0} + e^{2\pi i 0.j_n}\ket{1} \Big) \nonumber \\
\otimes &\Big( \ket{0} + e^{2\pi i 0.j_{n-1}j_n}\ket{1} \Big) \nonumber \\
&\ \ \vdots \nonumber \\
\otimes &\Big( \ket{0} + e^{2\pi i 0.j_1 j_2 \dots j_n}\ket{1} \Big)
\label{eq:qft_product_representation}
\end{align}

\subsection{การสังเคราะห์วงจรควอนตัมและความซับซ้อน}
จากสมการ (\ref{eq:qft_product_representation}) วงจร QFT สามารถสร้างขึ้นได้อย่างเป็นระบบโดยใช้เพียงสองประเภทเกตเท่านั้น
\begin{enumerate}
    \item เกตฮาดามาร์ด ($H$) เพื่อสร้างการซ้อนทับและเฟสจากบิต $j_k$ ตัวแรก
    \item เกตเปลี่ยนเฟสควบคุมสองคิวบิต (Controlled-Phase Rotation Gate ---- $R_k$)
    \begin{equation}
    R_k = \begin{pmatrix} 1 & 0 \\ 0 & e^{2\pi i / 2^k} \end{pmatrix}
    \end{equation}
\end{enumerate}

ในการแปลงคิวบิตที่ 1 ต้องใช้เกต $H$ จำนวน 1 ตัว และเกตควบคุม $R_2, R_3, \dots, R_n$ จำนวน $n-1$ ตัว สำหรับคิวบิตที่สองต้องใช้เกต $1 + (n-2)$ ตัว ทำเช่นนี้ไปเรื่อยๆ จนถึงคิวบิตสุดท้าย จำนวนเกตรวมทั้งหมดที่ต้องใช้ในวงจร QFT คือ
\begin{equation}
N_{\text{gates}} = \sum_{m=1}^n m = \frac{n(n+1)}{2} = O(n^2)
\label{eq:qft_gate_complexity}
\end{equation}
เมื่อเปรียบเทียบกับขั้นตอนวิธีการแปลงฟูเรียร์แบบเร็วบนคอมพิวเตอร์คลาสสิก (Fast Fourier Transform ---- FFT) ซึ่งต้องใช้การคำนวณ $O(N \log N) = O(n 2^n)$ การดำเนินการ จะเห็นได้ว่า QFT มีความซับซ้อนเพียง $O(n^2)$ ซึ่งให้การเพิ่มความเร็วแบบทวีคูณเอกซ์โพเนนเชียล (Exponential Speedup) เหนือ FFT คลาสสิกอย่างมหาศาล

% ==========================================================
% 10.5 การประเมินเฟสควอนตัมและอัลกอริทึมของชอร์
% ==========================================================
\section{การประเมินเฟสควอนตัมและอัลกอริทึมของชอร์}
\index{การประเมินเฟสควอนตัม (Quantum Phase Estimation ---- QPE)}
\index{อัลกอริทึมของชอร์ (Shor's Algorithm)}
\index{การแยกตัวประกอบเฉพาะ (Prime Factorization)}
\index{การถอดรหัสลับ RSA (RSA Cryptanalysis)}
\index{รหัสลับยุคหลังควอนตัม (Post-Quantum Cryptography ---- PQC)}

\subsection{การประเมินเฟสควอนตัม (Quantum Phase Estimation ---- QPE)}
โพรโทคอลที่เป็นกระดูกสันหลังของอัลกอริทึมควอนตัมขั้นสูงเกือบทั้งหมดคือ \textbf{การประเมินเฟสควอนตัม (QPE)} กำหนดให้ $U$ เป็นตัวดำเนินการเอกภาพ และ $\ket{u}$ เป็นสถานะไอเกนของ $U$ โดยมีค่าไอเกนเป็น $e^{2\pi i \theta}$
\begin{equation}
U \ket{u} = e^{2\pi i \theta} \ket{u} \qquad (0 \le \theta < 1)
\end{equation}
เป้าหมายของ QPE คือการประมาณค่าเฟส $\theta$ ให้มีความแม่นยำ $t$ บิต วงจร QPE ประกอบด้วยการเตรียม $t$ คิวบิตนับในสถานะ $\ket{0}$ และรีจิสเตอร์เป้าหมายในสถานะ $\ket{u}$ จากนั้นประยุกต์ใช้เกตควบคุมกำลังสอง $C-U^{2^k}$ แล้วปิดท้ายด้วยการแปลงฟูเรียร์ควอนตัมผกผัน ($\text{QFT}^\dagger$) ส่งผลให้รีจิสเตอร์นับยุบตัวลงสู่ค่าฐานสองของ $\theta$ ด้วยความน่าจะเป็นสูงยิ่ง

\subsection{อัลกอริทึมการแยกตัวประกอบของชอร์ (Shor's Algorithm)}
ในคริสต์ศักราช 1994 ปีเตอร์ ชอร์ (Peter Shor) แห่งเบลล์แล็บส์ ได้สร้างแรงสั่นสะเทือนครั้งประวัติศาสตร์แก่วงการความมั่นคงปลอดภัยไซเบอร์โลก โดยการเสนออัลกอริทึมควอนตัมที่สามารถแยกตัวประกอบของจำนวนเต็มขนาดใหญ่ $N = p \cdot q$ ได้ภายในเวลาพหุนาม
\begin{equation}
T_{\text{Shor}} = O\big( (\log N)^3 \big)
\label{eq:shor_time_complexity}
\end{equation}
ในขณะที่ขั้นตอนวิธีคลาสสิกที่ดีที่สุดในปัจจุบันคือ General Number Field Sieve (GNFS) ต้องใช้เวลาแบบกึ่งเอกซ์โพเนนเชียล $O\left(\exp\left(c \sqrt[3]{\log N (\log \log N)^2}\right)\right)$

ขั้นตอนการทำงานของอัลกอริทึมชอร์อาศัยการแปลงปัญหาการแยกตัวประกอบไปสู่ \textbf{ปัญหาการหาคาบของฟังก์ชันเลขชี้กำลังมอดุโล (Order-Finding Problem)} ดังนี้
\begin{enumerate}
    \item สุ่มเลือกจำนวนเต็ม $a$ โดยที่ $1 < a < N$ ตรวจสอบว่า $\gcd(a, N) = 1$ หรือไม่ หากไม่เท่ากับ 1 แสดงว่าพบตัวประกอบแล้ว
    \item ใช้คอมพิวเตอร์ควอนตัมหาค่าคาบ $r$ (ลำดับต่ำสุด) ที่ทำให้
    \begin{equation}
    a^r \equiv 1 \pmod N
    \label{eq:order_finding_modular}
    \end{equation}
    โดยการสร้างฟังก์ชัน $f(x) = a^x \bmod N$ แล้วใช้กระบวนการ QPE ร่วมกับ $\text{QFT}^\dagger$ ในการสกัดคาบ $r$
    \item หากคาบ $r$ ที่ได้เป็นเลขคู่ เราสามารถแยกตัวประกอบทางพีชคณิตได้เป็น
    \begin{equation}
    a^r - 1 = (a^{r/2} - 1)(a^{r/2} + 1) \equiv 0 \pmod N
    \end{equation}
    \item คำนวณตัวหารร่วมมาก $\gcd(a^{r/2} - 1, N)$ และ $\gcd(a^{r/2} + 1, N)$ โดยใช้อัลกอริทึมแบบยุคลิดบนคอมพิวเตอร์คลาสสิก ซึ่งจะให้ตัวประกอบเฉพาะ $p$ และ $q$ ในเวลาอันรวดเร็ว
\end{enumerate}

\subsection{ผลกระทบต่อวิทยาการรหัสลับและมาตรฐาน PQC}
ระบบความปลอดภัยของธุรกรรมการเงินและโครงสร้างพื้นฐานอินเทอร์เน็ตเกือบทั้งหมดในปัจจุบันพึ่งพาอัลกอริทึมการเข้ารหัสกุญแจสาธารณะ RSA และ ECC (Elliptic Curve Cryptography) ซึ่งตั้งอยู่บนสมมติฐานว่าการแยกตัวประกอบจำนวนเฉพาะขนาด 2048 บิต หรือการหาลอการิทึมไม่ต่อเนื่อง ต้องใช้เวลาคำนวณหลายพันล้านปี การมาถึงของอัลกอริทึมชอร์ทำให้ระบบเหล่านี้สามารถถูกถอดรหัสได้ภายในไม่กี่ชั่วโมง ส่งผลให้สถาบันมาตรฐานและเทคโนโลยีแห่งชาติสหรัฐอเมริกา (NIST) ต้องประกาศมาตรฐาน \textbf{วิทยาการรหัสลับยุคหลังควอนตัม (Post-Quantum Cryptography ---- PQC)} ในปี ค.ศ. 2024 เช่น ML-KEM (Kyber) และ ML-DSA (Dilithium) ซึ่งใช้โครงสร้างแลตทิซ (Lattice-based Cryptography) ที่ทนทานต่อการโจมตีด้วยคอมพิวเตอร์ควอนตัม

% ==========================================================
% กล่องปฏิบัติการ WebAR & Image Tracking: แล็บที่ 11
% ==========================================================
\begin{qrcodelabbox}{การแปลงฟูเรียร์ควอนตัมและอัลกอริทึมของชอร์ WebAR}{https://tsanaphysics.github.io/ar-Quantum/lab11_qft_shor_algorithm.html}
\begin{minipage}[c]{0.70\textwidth}
\textbf{คำแนะนำการทดลองเสมือนจริงสามมิติ (Interactive QFT \& Shor's Algorithm Lab)}
\begin{enumerate}
    \item สแกนรหัสตอบสนองเร็ว (QR Code) ด้านขวาเพื่อเปิดห้องปฏิบัติการเสมือนจริง WebAR บนเบราว์เซอร์
    \item ส่องกล้องหรือเปิดโหมด 3D Interactive เพื่อสำรวจ \textbf{โครงสร้างวงจรการแปลงฟูเรียร์ควอนตัม (QFT 3-Qubit Circuit)} ที่หมุนและวิเคราะห์ได้แบบ 360 องศา
    \item ทดลองกำหนดจำนวนเต็ม $N$ (เช่น $N = 15$ หรือ $N = 21$) และค่าฐานสุ่ม $a$ เพื่อสังเกตการณ์การสกัดคาบของฟังก์ชันเลขชี้กำลังมอดุโล $f(x) = a^x \bmod N$ ผ่านการแทรกสอดควอนตัม
    \item สังเกตการณ์การยุบตัวของสถานะในรีจิสเตอร์นับ และการคำนวณตัวหารร่วมมาก $\gcd(a^{r/2} \pm 1, N)$ เพื่อแยกตัวประกอบเฉพาะออกมาเป็นขั้นตอนอย่างชัดเจน
    \item ควบคุมการหมุนวงจรและเลือกขั้นตอนการทำงานได้ด้วยการเคลื่อนไหวมือผ่านกล้องหน้า (AI Touchless Gesture)
\end{enumerate}
\end{minipage}
\hfill
\begin{minipage}[c]{0.25\textwidth}
\centering
\qrcode[height=2.8cm]{https://tsanaphysics.github.io/ar-Quantum/lab11_qft_shor_algorithm.html}
\vspace{4pt}
\scriptsize\textbf{สแกนเปิด LAB 11 WebAR}
\end{minipage}
\end{qrcodelabbox}

% ==========================================================
% 10.6 ขั้นตอนวิธีแบบผสมผสานในยุค NISQ (VQE และ QAOA)
% ==========================================================
\section{ขั้นตอนวิธีแบบผสมผสานคลาสสิก-ควอนตัมในยุค NISQ}
\index{ยุค NISQ (Noisy Intermediate-Scale Quantum)}
\index{อัลกอริทึม VQE (Variational Quantum Eigensolver)}
\index{อัลกอริทึม QAOA (Quantum Approximate Optimization)}
\index{เคมีควอนตัม (Quantum Chemistry)}

แม้ว่าอัลกอริทึมของชอร์จะมีพลังมหาศาล แต่การนำไปใช้งานจริงกับกุญแจ RSA-2048 จำเป็นต้องใช้คิวบิตที่ทนทานต่อความผิดพร่อง (Fault-Tolerant Logical Qubits) หลายพันคิวบิต ซึ่งต้องการคิวบิตทางกายภาพนับล้านตัว จอห์น เพรสคิลล์ (John Preskill) แห่งสถาบัน Caltech จึงได้บัญญัตินามยุคปัจจุบันของคอมพิวเตอร์ควอนตัมว่า \textbf{ยุค NISQ (Noisy Intermediate-Scale Quantum)} ซึ่งมีจำนวนคิวบิตประมาณ 50 ถึง 1,000 คิวบิต และยังไม่มีระบบแก้ความผิดพลาดที่สมบูรณ์

ในยุค NISQ ขั้นตอนวิธีที่ทรงประสิทธิภาพที่สุดคือ \textbf{อัลกอริทึมแบบผสมผสานคลาสสิก-ควอนตัม (Hybrid Classical-Quantum Algorithms)}

\subsection{อัลกอริทึมแปรผันไอเกนควอนตัม (VQE)}
เสนอโดย อัลแบร์โต เปรูซโซ (Alberto Peruzzo) และคณะในปี ค.ศ. 1914 \textbf{VQE (Variational Quantum Eigensolver)} ถูกออกแบบมาเพื่อคำนวณหาพลังงานสถานะพื้น (Ground State Energy) ของโมเลกุลในเคมีควอนตัมและฟิสิกส์สสารควบแน่น โดยอาศัย \textbf{ระเบียบวิธีแปรผันของเรย์ลีห์-ริทซ์ (Rayleigh-Ritz Variational Principle)}
\begin{equation}
E_0 \le \langle \hat{H} \rangle_{\boldsymbol{\theta}} = \frac{\bra{\psi(\boldsymbol{\theta})} \hat{H} \ket{\psi(\boldsymbol{\theta})}}{\braket{\psi(\boldsymbol{\theta})|\psi(\boldsymbol{\theta})}}
\label{eq:vqe_variational_bound}
\end{equation}

ขั้นตอนการทำงานของ VQE เป็นวงรอบปิด (Closed Loop) ดังนี้
\begin{enumerate}
    \item \textbf{ส่วนประมวลผลควอนตัม (QPU)} ---- เตรียมสถานะควอนตัมที่ขึ้นกับพารามิเตอร์ $\boldsymbol{\theta} = (\theta_1, \theta_2, \dots, \theta_p)$ ผ่านวงจรที่เรียกว่า \textbf{อันซัทซ์ (Ansatz Circuit)} ---- $\ket{\psi(\boldsymbol{\theta})} = U(\boldsymbol{\theta})\ket{0}$
    \item \textbf{การวัดค่าคาดหวัง} ---- ฮามิลโทเนียนของโมเลกุลถูกเขียนกระจายในรูปของผลรวมเทนเซอร์ของเมทริกซ์เพาลี $\hat{H} = \sum_i c_i \hat{P}_i$ QPU จะทำการวัดค่าคาดหวัง $\langle \hat{P}_i \rangle$ แต่ละพจน์ แล้วส่งค่าพลังงานรวม $E(\boldsymbol{\theta}) = \sum_i c_i \langle \hat{P}_i \rangle_{\boldsymbol{\theta}}$ ไปยังคอมพิวเตอร์คลาสสิก
    \item \textbf{ส่วนปรับปรุงคลาสสิก (CPU/GPU Optimizer)} ---- คอมพิวเตอร์คลาสสิกใช้อัลกอริทึมการหาค่าต่ำสุด (เช่น COBYLA, SPSA หรือ Gradient Descent) เพื่อคำนวณชุดพารามิเตอร์ใหม่ $\boldsymbol{\theta}_{\text{new}}$ แล้วส่งกลับไปอัปเดตมุมของเกตใน QPU
\end{enumerate}
กระบวนการนี้จะวนซ้ำจนกระทั่งพลังงานลู่เข้าสู่ค่าต่ำสุด ซึ่งจะประมาณพลังงานสถานะพื้นของโมเลกุลได้อย่างแม่นยำสูง

\subsection{อัลกอริทึมการหาค่าเหมาะสมที่สุดโดยประมาณ (QAOA)}
เสนอโดย เอ็ดเวิร์ด ฟาร์ฮี (Edward Farhi), เจฟฟรีย์ โกลด์สโตน (Jeffrey Goldstone) และ แซม กัตมันน์ (Sam Gutmann) ในปี ค.ศ. 2014 \textbf{QAOA (Quantum Approximate Optimization Algorithm)} เป็นอัลกอริทึมที่ออกแบบมาเพื่อแก้ปัญหาการหาค่าเหมาะสมที่สุดเชิงการจัด (Combinatorial Optimization) เช่น ปัญหา Max-Cut ปัญหาการจัดเส้นทางโลจิสติกส์ และปัญหาการวางแผนโครงข่ายไฟฟ้า

วงจร QAOA ใช้การสลับชั้นของสองฮามิลโทเนียนเป็นจำนวน $p$ ชั้น
\begin{equation}
\ket{\gamma, \beta} = \prod_{k=1}^p \left( e^{-i\beta_k \hat{H}_M} e^{-i\gamma_k \hat{H}_C} \right) \ket{s}
\label{eq:qaoa_state_ansatz}
\end{equation}
โดยที่ $\hat{H}_C$ คือฮามิลโทเนียนต้นทุน (Cost Hamiltonian) ที่เข้ารหัสเงื่อนไขของปัญหา และ $\hat{H}_M = \sum_i \sigma_{x,i}$ คือฮามิลโทเนียนตัวผสม (Mixer Hamiltonian) คอมพิวเตอร์คลาสสิกจะปรับค่าพารามิเตอร์ $(\boldsymbol{\gamma}, \boldsymbol{\beta})$ เพื่อให้ได้คำตอบที่มีอัตราส่วนการประมาณ (Approximation Ratio) สูงที่สุด

% ==========================================================
% 10.7 รหัสแก้ความผิดพลาดควอนตัม (QEC)
% ==========================================================
\section{รหัสแก้ความผิดพลาดควอนตัมและเซอร์เฟสโค้ด}
\index{รหัสแก้ความผิดพลาดควอนตัม (Quantum Error Correction)}
\index{เซอร์เฟสโค้ด (Surface Codes)}
\index{คิวบิตเชิงตรรกะ (Logical Qubits)}
\index{ตัวดำเนินการสเตบิไลเซอร์ (Stabilizer Formalism)}

เนื่องจากสถานะควอนตัมมีความเปราะบางอย่างยิ่งยวดต่อการรบกวนจากสิ่งแวดล้อม (Decoherence) การสร้างคอมพิวเตอร์ควอนตัมขนาดใหญ่ที่ทำงานได้อย่างสมบูรณ์แบบจึงจำเป็นต้องพึ่งพา \textbf{รหัสแก้ความผิดพลาดควอนตัม (Quantum Error Correction หรือ QEC)}

ทว่าการแก้ความผิดพลาดในระบบควอนตัมต้องเผชิญกับอุปสรรคพื้นฐานสามประการที่ไม่มีในระบบคลาสสิก
\begin{enumerate}
    \item \textbf{ทฤษฎีบทห้ามโคลน} ---- เราไม่สามารถทำสำเนาคิวบิตเพื่อสำรองข้อมูลแบบรหัสความซ้ำซ้อนคลาสสิก ($0 \to 000$) ได้
    \item \textbf{ความผิดพลาดแบบต่อเนื่อง} ---- ความผิดพลาดในคิวบิตไม่ได้มีเพียงการกลับบิต ($0 \leftrightarrow 1$) แต่รวมถึงการหมุนเฟสของสถานะเป็นมุมต่อเนื่องใดๆ
    \item \textbf{การวัดทำลายการซ้อนทับ} ---- การตรวจวัดสถานะเพื่อตรวจสอบความผิดพลาดโดยตรงจะทำให้ฟังก์ชันคลื่นยุบตัวและสูญเสียสารสนเทศควอนตัม
\end{enumerate}

ปีเตอร์ ชอร์ ได้แก้ปัญหาอันซับซ้อนนี้ในปี ค.ศ. 1995 โดยการเสนอ \textbf{รหัส 9 คิวบิตของชอร์ (Shor 9-Qubit Code)} ซึ่งแสดงให้เห็นว่า เราสามารถตรวจจับความผิดพลาดได้โดยการวัดเฉพาะ \textbf{กลุ่มอาการความผิดพลาด (Error Syndrome)} ผ่านตัวดำเนินการสเตบิไลเซอร์ โดยไม่ต้องวัดสถานะของคิวบิตข้อมูลโดยตรง

\subsection{เซอร์เฟสโค้ด (Surface Codes)}
ในปัจจุบัน สถาปัตยกรรม QEC ที่ได้รับการยอมรับอย่างกว้างขวางที่สุดคือ \textbf{เซอร์เฟสโค้ด (Surface Codes)} ซึ่งเสนอโดย อเล็กซี คิเทียฟ (Alexei Kitaev) ในปี ค.ศ. 1997 คิวบิตกายภาพจะถูกจัดเรียงเป็นโครงตาข่ายสองมิติบนระนาบ โดยแบ่งออกเป็น
\begin{itemize}
    \item \textbf{คิวบิตข้อมูล (Data Qubits)} ทำหน้าที่เก็บสารสนเทศควอนตัมเชิงตรรกะ
    \item \textbf{คิวบิตช่วยตรวจจับ (Syndrome Ancilla Qubits)} ทำหน้าที่วัดตัวดำเนินการสเตบิไลเซอร์ชนิด $Z$-Plaquette (ตรวจจับความผิดพลาดแบบ Bit-flip ---- $X$) และชนิด $X$-Star (ตรวจจับความผิดพลาดแบบ Phase-flip ---- $Z$)
\end{itemize}

จุดเด่นที่สุดของเซอร์เฟสโค้ดคือ \textbf{ต้องการเพียงการเชื่อมต่อระหว่างคิวบิตข้างเคียงที่ติดกันเท่านั้น (Nearest-Neighbor Coupling)} และมีเกณฑ์ความผิดพลาดวิกฤต (Fault-Tolerant Threshold) สูงถึงประมาณร้อยละ $1$ ($p_{\text{th}} \approx 10^{-2}$) ซึ่งสอดคล้องกับขีดความสามารถของฮาร์ดแวร์ในปัจจุบัน

ในปี ค.ศ. 2023 ทีมวิจัย Google Quantum AI ได้แสดงหลักไมล์สำคัญเป็นครั้งแรกของโลก โดยพิสูจน์ว่าเมื่อเพิ่มขนาดของเซอร์เฟสโค้ดจากระยะ $d=3$ (คิวบิตกายภาพ 17 ตัว) สู่ระยะ $d=5$ (คิวบิตกายภาพ 49 ตัว) อัตราความผิดพลาดของคิวบิตเชิงตรรกะลดลงอย่างชัดเจน ซึ่งเป็นข้อพิสูจน์เชิงประจักษ์ว่าคอมพิวเตอร์ควอนตัมทนทานต่อความผิดพร่องสามารถสร้างขึ้นได้จริงตามทฤษฎี

% ==========================================================
% 10.8 สถาปัตยกรรมฮาร์ดแวร์ควอนตัมยุคใหม่
% ==========================================================
\section{สถาปัตยกรรมฮาร์ดแวร์ควอนตัมคอมพิวติง}
\index{คิวบิตทรานสมอน (Transmon Qubits)}
\index{ไอออนในกับดัก (Trapped Ions)}
\index{อะตอมเป็นกลาง (Neutral Atoms)}
\index{ควอนตัมโฟตอนิกส์ (Photonic Quantum Computing)}

การแข่งขันทางเทคโนโลยีในการสร้างคอมพิวเตอร์ควอนตัมกำลังดำเนินไปอย่างเข้มข้นผ่านแพลตฟอร์มฮาร์ดแวร์ทางฟิสิกส์ที่หลากหลาย ดังแสดงการเปรียบเทียบในตารางที่ \ref{tab:quantum_hardware_comparison}

\begin{table}[H]
\centering
\caption{การเปรียบเทียบคุณสมบัติทางฟิสิกส์ของแพลตฟอร์มฮาร์ดแวร์ควอนตัมคอมพิวติงชั้นนำ}
\label{tab:quantum_hardware_comparison}
\begin{tabular}{lcccc}
\toprule
\textbf{แพลตฟอร์ม} & \textbf{เวลาเชื่อมแน่น ($T_2$)} & \textbf{ความเร็วเกต} & \textbf{ความแม่นยำ 2 คิวบิต} & \textbf{อุณหภูมิทำงาน} \\
\midrule
ตัวนำยิ่งยวด (Transmon) & $50 - 300\ \mu\text{s}$ & $10 - 50\text{ ns}$ & $> 99.5\%$ & $\sim 15\text{ mK}$ \\
ไอออนในกับดัก (Trapped Ion) & วินาที ถึง นาที & $10 - 100\ \mu\text{s}$ & $> 99.9\%$ & อุณหภูมิห้อง / $4\text{ K}$ \\
อะตอมเป็นกลาง (Rydberg) & มิลลิวินาที ถึง วินาที & $0.1 - 1\ \mu\text{s}$ & $> 99.5\%$ & $\sim 10\ \mu\text{K}$ \\
โฟตอนิกส์ (Photonic) & ไม่จำกัด (ที่ความเร็วแสง) & ระดับพิโกวินาที & $\sim 99\%$ & อุณหภูมิห้อง \\
สปินในซิลิคอน (Si Spin) & มิลลิวินาที & $10 - 100\text{ ns}$ & $> 99.5\%$ & $\sim 1\text{ K}$ \\
\bottomrule
\end{tabular}
\end{table}

\begin{enumerate}
    \item \textbf{วงจรตัวนำยิ่งยวดทรานสมอน (Superconducting Transmon)} ---- ผู้นำในปัจจุบันโดย IBM, Google และ Rigetti อาศัยรอยต่อโจเซฟสัน (Josephson Junction) สร้างความไม่เป็นฮาร์มอนิกของระดับพลังงาน ข้อดีคือความเร็วเกตสูงมากในระดับสิบนาโนวินาที และผลิตได้ด้วยเทคโนโลยีลิโธกราฟีของเซมิคอนดักเตอร์
    \item \textbf{ไอออนในกับดักสนามแม่เหล็กไฟฟ้า (Trapped Ions)} ---- พัฒนาโดย Quantinuum และ IonQ ใช้ไอออนของธาตุอัลคาไลน์เอิร์ธ (เช่น $^{171}\text{Yb}^+$) ถูกกักขังด้วยกับดักเพาลี (Paul Trap) มีความเชื่อมแน่นยาวนานที่สุดและคิวบิตทุกตัวสามารถเชื่อมต่อถึงกันได้ทั้งหมด (All-to-All Connectivity)
    \item \textbf{อะตอมเป็นกลางในคีมจับแสง (Neutral Atoms in Optical Tweezers)} ---- พัฒนาโดย QuEra และมหาวิทยาลัยฮาร์วาร์ด กักขังอะตอมเดี่ยว (เช่น รูบิเดียมหรือซีเซียม) ด้วยลำแสงเลเซอร์นับพันจุด และสร้างอันตรกิริยาระหว่างคิวบิตผ่านการกระตุ้นสู่สถานะริดเบิร์ก (Rydberg Blockade) ซึ่งสามารถขยายขนาดสู่หลายพันคิวบิตได้อย่างรวดเร็ว
\end{enumerate}

% ==========================================================
% 10.9 เซนเซอร์ควอนตัมและมาตรวิทยาควอนตัม
% ==========================================================
\section{เซนเซอร์ควอนตัมและมาตรวิทยาควอนตัมในยุคใหม่}
\index{เซนเซอร์ควอนตัม (Quantum Sensing)}
\index{จุดบกพร่องไนโตรเจน-ช่องว่าง (NV Center)}
\index{ขีดจำกัดไฮเซนเบิร์ก (Heisenberg Limit)}
\index{เกษตรดิจิทัลแม่นยำสูง}

นอกจากพลังในการคำนวณ เทคโนโลยีควอนตัมยุคที่สองยังได้ปฏิวัติวิทยาการการวัดและการรับรู้ด้วย \textbf{เซนเซอร์ควอนตัม (Quantum Sensors)} ซึ่งอาศัยความไวขั้นสูงสุดของสถานะซ้อนทับและความพัวพันต่อสนามแม่เหล็ก สนามไฟฟ้า อุณหภูมิ และความโน้มถ่วงภายนอก

ตามหลักการของมาตรวิทยาควอนตัม (Quantum Metrology) หากใช้ระบบคลาสสิกที่ไม่มีความพัวพันจำนวน $N$ อนุภาค ความไม่แน่นอนในการวัดจะถูกจำกัดด้วย \textbf{ขีดจำกัดมาตรฐานควอนตัม (Standard Quantum Limit ---- SQL)}
\begin{equation}
\Delta \theta_{\text{SQL}} \ge \frac{1}{\sqrt{N}}
\label{eq:sql_limit}
\end{equation}
ทว่าหากนำอนุภาคทั้ง $N$ ตัวมาพัวพันกันในสถานะพิเศษ (เช่น สถานะ NOON หรือสถานะบีบอัด Squeezed State) ความแม่นยำจะทะลุขอบเขตคลาสสิกเข้าสู่ \textbf{ขีดจำกัดของไฮเซนเบิร์ก (Heisenberg Limit)}
\begin{equation}
\Delta \theta_{\text{Heisenberg}} \ge \frac{1}{N}
\label{eq:heisenberg_limit}
\end{equation}
ซึ่งให้ความไวสูงกว่าเดิมถึง $\sqrt{N}$ เท่า

\subsection{จุดบกพร่องไนโตรเจน-ช่องว่างในเพชร (NV Centers in Diamond)}
ระบบเซนเซอร์ควอนตัมที่มีความโดดเด่นสูงสุดสำหรับการประยุกต์ใช้งานจริงคือ \textbf{จุดบกพร่องไนโตรเจน-ช่องว่าง (Nitrogen-Vacancy Center หรือ NV Center)} ซึ่งเป็นข้อบกพร่องระดับอะตอมในโครงผลึกเพชร โดยมีอะตอมไนโตรเจนเข้ามาแทนที่คาร์บอนและมีช่องว่างข้างเคียง ระบบนี้ทำหน้าที่เป็นสปินอิเล็กตรอนเดี่ยว ($S=1$) ที่สามารถเตรียมสถานะ ควบคุม และอ่านค่าได้ด้วยแสงเลเซอร์และการสั่นพ้องแม่เหล็กเชิงแสง (Optically Detected Magnetic Resonance ---- ODMR) ที่อุณหภูมิห้อง

การประยุกต์ใช้ในการเกษตรดิจิทัลและสิ่งแวดล้อม
\begin{itemize}
    \item \textbf{การตรวจวัดไอออนธาตุอาหารในพืชแบบเรียลไทม์} ---- เข็มวัดระดับนาโนเมตรที่ทำจากเพชร NV Center สามารถแทงเข้าสู่เซลล์ท่อน้ำเลี้ยง (Xylem/Phloem) ของพืชเศรษฐกิจ เพื่อตรวจวัดความเข้มข้นของไอออนฟอสเฟต โพแทสเซียม และไนเตรต โดยไม่ทำลายเนื้อเยื่อพืช
    \item \textbf{การวิเคราะห์จุลชีววิทยาและฟิสิกส์ของดิน} ---- ตรวจวัดการเปลี่ยนแปลงของความร้อนระดับไมโครเคลวินและปฏิกิริยาเคมีรีดอกซ์ในดินชั้นลึก เพื่อทำนายการทำงานของจุลินทรีย์ตรึงไนโตรเจน
    \item \textbf{การวัดความโน้มถ่วงยิ่งยวด (Quantum Gravimetry)} ---- ใช้อินเทอร์เฟอโรเมทรีของอะตอมเย็น (Cold Atom Gravimeter) วัดความเร่งโน้มถ่วง $g$ ด้วยความละเอียด $10^{-9} g$ เพื่อตรวจจับการเปลี่ยนแปลงของตาน้ำใต้ดินและการกักเก็บคาร์บอนในแปลงเกษตรขนาดใหญ่
\end{itemize}

% ==========================================================
% 10.10 การจำลองด้วยภาษา Python และ Qiskit
% ==========================================================
\section{การจำลองวงจรโกรเวอร์และ VQE ด้วย Python และ Qiskit}

\begin{tcolorbox}[colback=codeBackground, colframe=darkNavy, arc=2mm, title=\textbf{สคริปต์ Python ที่ 10.1 ---- การจำลองอัลกอริทึมของโกรเวอร์ขนาด 3 คิวบิตด้วย Qiskit}]
\begin{lstlisting}[language=Python]
# ==========================================================
# 10_grover_3qubit_simulation.py
# การจำลองอัลกอริทึมของโกรเวอร์สำหรับค้นหาสถานะ |101> (3 คิวบิต)
# สอดคล้องกับระเบียบวิธีของ Dr. Chewa Thassana
# ==========================================================
import numpy as np
from qiskit import QuantumCircuit
from qiskit.quantum_info import Statevector

def create_grover_circuit():
    # กำหนดวงจรขนาด 3 คิวบิต (N = 2^3 = 8 รายการ)
    qc = QuantumCircuit(3)
    
    # 1. เตรียมสถานะซ้อนทับสม่ำเสมอ |s>
    qc.h([0, 1, 2])
    qc.barrier()
    
    # 2. ออราเคิลสำหรับเป้าหมาย |w> = |101> (q2=1, q1=0, q0=1)
    # พลิกเฟสเฉพาะสถานะ |101> ด้วยการกลับค่าบิต 1 แล้วใช้ Multi-Controlled-Z
    qc.x(1)
    qc.h(2)
    qc.ccx(0, 1, 2)
    qc.h(2)
    qc.x(1)
    qc.barrier()
    
    # 3. ตัวดำเนินการแพร่กระจาย (Grover Diffuser ---- 2|s><s| - I)
    qc.h([0, 1, 2])
    qc.x([0, 1, 2])
    qc.h(2)
    qc.ccx(0, 1, 2)
    qc.h(2)
    qc.x([0, 1, 2])
    qc.h([0, 1, 2])
    
    return qc

if __name__ == "__main__":
    qc = create_grover_circuit()
    sv = Statevector.from_instruction(qc)
    probs = sv.probabilities_dict()
    
    print("ผลลัพธ์การจำลองการค้นหาของโกรเวอร์ 1 รอบ (N=8, M=1) ----")
    for state, p in sorted(probs.items()):
        bar = "#" * int(p * 50)
        print(f"|{state}> ---- ความน่าจะเป็น = {p:.4f} ({p*100:5.1f}%) {bar}")
    
    p_target = probs.get("101", 0.0)
    print(f"\nความน่าจะเป็นของสถานะเป้าหมาย |101> พุ่งสูงถึง {p_target*100:.2f}%")
\end{lstlisting}
\end{tcolorbox}

\begin{tcolorbox}[colback=codeBackground, colframe=emeraldGlow, arc=2mm, title=\textbf{สคริปต์ Python ที่ 10.2 ---- การจำลองขั้นตอนวิธี VQE สำหรับ 2 คิวบิตด้วย NumPy}]
\begin{lstlisting}[language=Python]
# ==========================================================
# 10_vqe_minimal_solver.py
# การจำลองขั้นตอนวิธี Variational Quantum Eigensolver (VQE)
# ==========================================================
import numpy as np
from scipy.optimize import minimize

# นิยามเมทริกซ์เพาลี
I = np.eye(2, dtype=complex)
X = np.array([[0, 1], [1, 0]], dtype=complex)
Z = np.array([[1, 0], [0, -1]], dtype=complex)

# ฮามิลโทเนียนจำลองของโมเลกุล ---- H = -1.05*(Z (x) I) - 1.05*(I (x) Z) + 0.4*(X (x) X)
H = -1.05 * np.kron(Z, I) - 1.05 * np.kron(I, Z) + 0.4 * np.kron(X, X)

# คำนวณค่าไอเกนต่ำสุดที่แท้จริงทางคณิตศาสตร์ (Ground State Energy)
exact_ground_energy = np.min(np.linalg.eigvalsh(H))

# นิยามฟังก์ชันสร้างสถานะแปรผัน (Ansatz) ---- Ry(theta1) (x) Ry(theta2) กระทำกับ |00>
def ansatz(params):
    t1, t2 = params
    ry1 = np.array([[np.cos(t1/2), -np.sin(t1/2)], [np.sin(t1/2), np.cos(t1/2)]])
    ry2 = np.array([[np.cos(t2/2), -np.sin(t2/2)], [np.sin(t2/2), np.cos(t2/2)]])
    psi0 = np.array([1, 0, 0, 0], dtype=complex)
    return np.kron(ry1, ry2) @ psi0

# ฟังก์ชันต้นทุนคำนวณค่าคาดหวังพลังงาน <psi(theta)|H|psi(theta)>
def energy_cost(params):
    psi = ansatz(params)
    return np.real(np.vdot(psi, H @ psi))

# ดำเนินการหาค่าต่ำสุดด้วยอัลกอริทึมคลาสสิก COBYLA
init_params = [0.5, 0.5]
res = minimize(energy_cost, init_params, method='COBYLA')

print("ผลการคำนวณ Variational Quantum Eigensolver (VQE) ----")
print(f"พลังงานสถานะพื้นจริง (Exact Ground Energy) = {exact_ground_energy:.6f} Hartree")
print(f"พลังงานที่คำนวณได้จาก VQE                  = {res.fun:.6f} Hartree")
print(f"ความคลาดเคลื่อนสัมบูรณ์                   = {abs(res.fun - exact_ground_energy):.6e}")
print(f"พารามิเตอร์มุมที่เหมาะสมที่สุด (theta1, theta2) = ({res.x[0]:.4f}, {res.x[1]:.4f})")
\end{lstlisting}
\end{tcolorbox}

% ==========================================================
% 10.11 ตัวอย่างโจทย์เชิงวิเคราะห์ตามระเบียบวิธี P-S-C-T
% ==========================================================
\newpage
\section{ตัวอย่างการแก้โจทย์เชิงวิเคราะห์ตามระเบียบวิธี P-S-C-T}

% ตัวอย่างที่ 10.1
\begin{psctbox}{ตัวอย่างที่ 10.1 ---- การวิเคราะห์เวกเตอร์สถานะในอัลกอริทึมดอยช์สำหรับ 1 คิวบิต}
\noindent\textbf{P - Problem Statement (โจทย์ปัญหา)} \\
กำหนดให้ออราเคิล 1 คิวบิตแทนฟังก์ชัน $f(x) = x$ (นั่นคือ $f(0) = 0$ และ $f(1) = 1$) จงแสดงการวิวัฒน์ของเวกเตอร์สถานะในทุกขั้นตอนของอัลกอริทึมดอยช์ และคำนวณความน่าจะเป็นของการวัดค่าคิวบิตแรก

\vspace{4pt}
\noindent\textbf{S - Strategy \& Principle (กลยุทธ์และหลักการ)} \\
เริ่มต้นจากสถานะ $\ket{\psi_0} = \ket{0}\ket{1}$ กระทำเกต $H \otimes H$ นำเข้าสู่ออราเคิล $U_f$ ซึ่งสร้างการสะท้อนเฟสตามสมการ $U_f\ket{x}\ket{-} = (-1)^{f(x)}\ket{x}\ket{-}$ แล้วกระทำเกต $H$ บนคิวบิตแรกก่อนการวัด

\vspace{4pt}
\noindent\textbf{C - Calculation \& Proof (การคำนวณและการพิสูจน์)} \\
ขั้นตอนที่ 1 ---- สถานะหลังผ่านเกตฮาดามาร์ดคู่
\begin{equation}
\ket{\psi_1} = (H \otimes H)\ket{01} = \left( \frac{\ket{0}+\ket{1}}{\sqrt{2}} \right) \otimes \left( \frac{\ket{0}-\ket{1}}{\sqrt{2}} \right) = \frac{\ket{0}\ket{-} + \ket{1}\ket{-}}{\sqrt{2}}
\end{equation}
ขั้นตอนที่ 2 ---- ผ่านเกตออราเคิล $U_f$ เนื่องจาก $f(0)=0 \implies (-1)^0 = +1$ และ $f(1)=1 \implies (-1)^1 = -1$
\begin{equation}
\ket{\psi_2} = U_f \ket{\psi_1} = \frac{(+1)\ket{0}\ket{-} + (-1)\ket{1}\ket{-}}{\sqrt{2}} = \left( \frac{\ket{0}-\ket{1}}{\sqrt{2}} \right) \otimes \ket{-} = \ket{-}\ket{-}
\end{equation}
ขั้นตอนที่ 3 ---- กระทำเกต $H$ บนคิวบิตแรก
\begin{equation}
\ket{\psi_3} = (H \otimes I)\ket{\psi_2} = (H\ket{-}) \otimes \ket{-} = \ket{1} \otimes \ket{-}
\end{equation}
เมื่อทำการวัดเฉพาะคิวบิตแรก ผลลัพธ์จะยุบตัวลงสู่สถานะ $\ket{1}$ ด้วยความน่าจะเป็นร้อยเปอร์เซ็นต์เต็ม ($P(1) = 1.0$)

\vspace{4pt}
\noindent\textbf{T - Takeaway \& Modern Context (ข้อคิดสรุปและบริบทสมัยใหม่)} \\
การวัดได้ $\ket{1}$ สรุปได้อย่างแน่นอนเด็ดขาดว่าฟังก์ชันเป็นแบบสมดุล (Balanced) ซึ่งสอดคล้องกับ $f(x)=x$ หากฟังก์ชันเป็นแบบคงตัว ผลลัพธ์ที่ได้จะเป็น $\ket{0}$ เสมอ การทดลองนี้แสดงให้เห็นว่าคอมพิวเตอร์ควอนตัมสามารถแยกแยะคุณสมบัติระดับสากล (Global Property) ของฟังก์ชันได้โดยการประเมินค่าเพียงครั้งเดียว
\end{psctbox}

% ตัวอย่างที่ 10.2
\begin{psctbox}{ตัวอย่างที่ 10.2 ---- การคำนวณรอบการทำซ้ำและความน่าจะเป็นของโกรเวอร์ในกรณีคำตอบเดี่ยว}
\noindent\textbf{P - Problem Statement (โจทย์ปัญหา)} \\
กำหนดฐานข้อมูลขนาด $N = 2^4 = 16$ รายการ โดยมีคำตอบเป้าหมายเดี่ยว $M = 1$ จงคำนวณมุมการหมุนของโกรเวอร์ $\theta$ จำนวนรอบการทำซ้ำที่เหมาะสม $R$ และความน่าจะเป็นของความสำเร็จ $P_{\text{success}}$

\vspace{4pt}
\noindent\textbf{S - Strategy \& Principle (กลยุทธ์และหลักการ)} \\
คำนวณมุมจาก $\sin(\theta/2) = \sqrt{M/N} = \sqrt{1/16} = 1/4$ จากนั้นหาจำนวนรอบเต็ม $R = \text{round}\left(\frac{\pi}{2\theta} - \frac{1}{2}\right)$ และคำนวณ $P_{\text{success}} = \sin^2\left(\frac{2R+1}{2}\theta\right)$

\vspace{4pt}
\noindent\textbf{C - Calculation \& Proof (การคำนวณและการพิสูจน์)} \\
คำนวณมุม $\theta$
\begin{equation}
\sin\left(\frac{\theta}{2}\right) = \frac{1}{4} = 0.25 \implies \frac{\theta}{2} = \arcsin(0.25) \approx 0.25268\text{ rad} \implies \theta \approx 0.50536\text{ rad}
\end{equation}
คำนวณจำนวนรอบการทำซ้ำที่เหมาะสมที่สุด
\begin{equation}
R = \text{round}\left( \frac{\pi/2}{0.50536} - 0.5 \right) = \text{round}(3.108 - 0.5) = \text{round}(2.608) = 3\text{ รอบ}
\end{equation}
คำนวณมุมรวมหลังทำซ้ำ $R=3$ รอบ
\begin{equation}
\phi_3 = \frac{2(3)+1}{2}\theta = \frac{7}{2}(0.50536) = 1.76876\text{ rad} \ (101.34^\circ)
\end{equation}
คำนวณความน่าจะเป็นของความสำเร็จ
\begin{equation}
P_{\text{success}} = \sin^2(1.76876) = (0.98048)^2 \approx 0.9613 = 96.13\%
\end{equation}

\vspace{4pt}
\noindent\textbf{T - Takeaway \& Modern Context (ข้อคิดสรุปและบริบทสมัยใหม่)} \\
การทำซ้ำเพียง 3 รอบสามารถเพิ่มความน่าจะเป็นของคำตอบจากเดิมเพียง $1/16 = 6.25\%$ ขึ้นสู่ $96.13\%$ หากทำซ้ำต่อไปเป็นรอบที่ 4 มุมรวมจะกลายเป็น $\phi_4 = 9/2(0.50536) = 2.274\text{ rad}$ ส่งผลให้ความน่าจะเป็นลดลงเหลือเพียง $58.1\%$ การเข้าใจพลวัตเชิงเรขาคณิตของการหมุนจึงมีความสำคัญอย่างยิ่งต่อการควบคุมอัลกอริทึม
\end{psctbox}

% ตัวอย่างที่ 10.3
\begin{psctbox}{ตัวอย่างที่ 10.3 ---- การแปลงฟูเรียร์ควอนตัมของสถานะฐาน 3 คิวบิต}
\noindent\textbf{P - Problem Statement (โจทย์ปัญหา)} \\
จงคำนวณการแปลงฟูเรียร์ควอนตัม (QFT) ของสถานะฐาน $\ket{4} \equiv \ket{100}_2$ ในระบบ 3 คิวบิต ($N=2^3=8$) และเขียนผลลัพธ์ในรูปผลคูณเทนเซอร์ของแต่ละคิวบิต

\vspace{4pt}
\noindent\textbf{S - Strategy \& Principle (กลยุทธ์และหลักการ)} \\
ใช้รูปแบบผลคูณเทนเซอร์ของ QFT จากสมการ (\ref{eq:qft_product_representation}) โดยที่ $j = 4 = 100_2$ ($j_1=1, j_2=0, j_3=0$) คำนวณเศษส่วนฐานสอง $0.j_3$, $0.j_2 j_3$, $0.j_1 j_2 j_3$

\vspace{4pt}
\noindent\textbf{C - Calculation \& Proof (การคำนวณและการพิสูจน์)} \\
พิจารณาเศษส่วนฐานสองในแต่ละพจน์
\begin{align}
0.j_3 &= 0.0_2 = 0 \implies e^{2\pi i (0)} = 1 \\
0.j_2 j_3 &= 0.00_2 = 0 \implies e^{2\pi i (0)} = 1 \\
0.j_1 j_2 j_3 &= 0.100_2 = \frac{1}{2} \implies e^{2\pi i (1/2)} = e^{i\pi} = -1
\end{align}
แทนค่าลงในรูปแบบผลคูณเทนเซอร์
\begin{align}
\text{QFT}\ket{100} &= \frac{1}{\sqrt{8}} \Big( \ket{0} + e^{2\pi i (0)}\ket{1} \Big) \otimes \Big( \ket{0} + e^{2\pi i (0)}\ket{1} \Big) \otimes \Big( \ket{0} + e^{2\pi i (1/2)}\ket{1} \Big) \nonumber \\
&= \frac{1}{\sqrt{8}} (\ket{0}+\ket{1}) \otimes (\ket{0}+\ket{1}) \otimes (\ket{0}-\ket{1}) \nonumber \\
&= \ket{+} \otimes \ket{+} \otimes \ket{-}
\end{align}

\vspace{4pt}
\noindent\textbf{T - Takeaway \& Modern Context (ข้อคิดสรุปและบริบทสมัยใหม่)} \\
สถานะผลลัพธ์คือ $\ket{++-}$ ซึ่งเป็นสถานะแยกกันได้อย่างสมบูรณ์แบบ การที่สถานะผลลัพธ์สามารถแยกตัวประกอบเทนเซอร์ได้ชี้ให้เห็นว่า QFT แปลงข้อมูลเชิงตำแหน่งไปเป็นข้อมูลเชิงเฟสของแต่ละคิวบิต ซึ่งเป็นคุณสมบัติหลักที่นำไปใช้ในอัลกอริทึมการประมาณเฟส (QPE)
\end{psctbox}

% ตัวอย่างที่ 10.4
\begin{psctbox}{ตัวอย่างที่ 10.4 ---- การแยกตัวประกอบ $N=15$ ด้วยขั้นตอนวิธีของชอร์}
\noindent\textbf{P - Problem Statement (โจทย์ปัญหา)} \\
ในการแยกตัวประกอบ $N=15$ ด้วยอัลกอริทึมของชอร์ หากสุ่มเลือกจำนวนเต็ม $a=7$ จงแสดงการหาคาบ $r$ ของฟังก์ชันมอดุโล $f(x) = 7^x \bmod 15$ และแสดงการแยกตัวประกอบหา $p$ และ $q$

\vspace{4pt}
\noindent\textbf{S - Strategy \& Principle (กลยุทธ์และหลักการ)} \\
คำนวณลำดับค่า $7^x \bmod 15$ สำหรับ $x = 1, 2, 3, \dots$ เพื่อหาค่าคาบ $r$ ที่ทำให้ $7^r \equiv 1 \pmod{15}$ จากนั้นตรวจสอบว่า $r$ เป็นเลขคู่หรือไม่ และคำนวณ $\gcd(7^{r/2} \pm 1, 15)$

\vspace{4pt}
\noindent\textbf{C - Calculation \& Proof (การคำนวณและการพิสูจน์)} \\
คำนวณค่าฟังก์ชันมอดุโล
\begin{align}
x = 1 &\implies 7^1 \bmod 15 = 7 \\
x = 2 &\implies 7^2 \bmod 15 = 49 \bmod 15 = 4 \\
x = 3 &\implies 7^3 \bmod 15 = (4 \times 7) \bmod 15 = 28 \bmod 15 = 13 \\
x = 4 &\implies 7^4 \bmod 15 = (13 \times 7) \bmod 15 = 91 \bmod 15 = 1
\end{align}
พบว่าคาบ $r = 4$ ซึ่งเป็นเลขคู่ \\
คำนวณ $a^{r/2} = 7^{4/2} = 7^2 = 49$ \\
คำนวณตัวประกอบเฉพาะผ่านฟังก์ชันตัวหารร่วมมาก (GCD)
\begin{align}
p &= \gcd(7^2 - 1, 15) = \gcd(48, 15) = 3 \\
q &= \gcd(7^2 + 1, 15) = \gcd(50, 15) = 5
\end{align}
ได้ตัวประกอบเฉพาะของ 15 คือ $3$ และ $5$ อย่างถูกต้องสมบูรณ์

\vspace{4pt}
\noindent\textbf{T - Takeaway \& Modern Context (ข้อคิดสรุปและบริบทสมัยใหม่)} \\
ในการทดลองประวัติศาสตร์เมื่อปี ค.ศ. 2001 ทีมนักวิจัย IBM Almaden ได้ใช้วงจรคิวบิตเรโซแนนซ์แม่เหล็กนิวเคลียร์ (NMR) ในการรันอัลกอริทึมชอร์แยกตัวประกอบ $N=15$ ได้สำเร็จเป็นครั้งแรกของโลก การลดรูปปัญหาเชิงพีชคณิตสู่การหาคาบแสดงให้เห็นความเชื่อมโยงอันลึกซึ้งระหว่างทฤษฎีจำนวนและกลศาสตร์ควอนตัม
\end{psctbox}

% ตัวอย่างที่ 10.5
\begin{psctbox}{ตัวอย่างที่ 10.5 ---- การคำนวณพลังงานสถานะพื้นของโมเลกุลด้วย VQE}
\noindent\textbf{P - Problem Statement (โจทย์ปัญหา)} \\
กำหนดฮามิลโทเนียนอย่างง่ายของระบบสองคิวบิต $\hat{H} = -Z_1 - Z_2 + 0.5 X_1 X_2$ หากใช้อันซัทซ์สถานะพัวพันพารามิเตอร์เดี่ยว $\ket{\psi(\theta)} = \cos(\theta)\ket{00} + \sin(\theta)\ket{11}$ จงหาฟังก์ชันค่าคาดหวังพลังงาน $E(\theta)$ และคำนวณมุม $\theta_{\text{min}}$ ที่ให้พลังงานสถานะพื้นต่ำสุด

\vspace{4pt}
\noindent\textbf{S - Strategy \& Principle (กลยุทธ์และหลักการ)} \\
คำนวณค่าคาดหวังของแต่ละตัวดำเนินการเพาลี $\langle Z_1 \rangle$, $\langle Z_2 \rangle$ และ $\langle X_1 X_2 \rangle$ บนสถานะ $\ket{\psi(\theta)}$ แล้วหาอนุพันธ์ $\frac{dE}{d\theta} = 0$

\vspace{4pt}
\noindent\textbf{C - Calculation \& Proof (การคำนวณและการพิสูจน์)} \\
พิจารณาการกระทำของตัวดำเนินการบนสถานะ $\ket{\psi(\theta)}$
\begin{align}
Z_1 \ket{\psi(\theta)} &= \cos(\theta)\ket{00} - \sin(\theta)\ket{11} \implies \bra{\psi}Z_1\ket{\psi} = \cos^2(\theta) - \sin^2(\theta) = \cos(2\theta) \\
Z_2 \ket{\psi(\theta)} &= \cos(\theta)\ket{00} - \sin(\theta)\ket{11} \implies \bra{\psi}Z_2\ket{\psi} = \cos(2\theta) \\
(X_1 X_2) \ket{\psi(\theta)} &= \cos(\theta)\ket{11} + \sin(\theta)\ket{00} \implies \bra{\psi}X_1 X_2\ket{\psi} = 2\cos(\theta)\sin(\theta) = \sin(2\theta)
\end{align}
รวมค่าคาดหวังพลังงาน
\begin{equation}
E(\theta) = -\cos(2\theta) - \cos(2\theta) + 0.5\sin(2\theta) = -2\cos(2\theta) + 0.5\sin(2\theta)
\end{equation}
หาจุดต่ำสุดโดยการดิฟเฟอเรนชิเอตเทียบกับ $\theta$
\begin{equation}
\frac{dE}{d\theta} = 4\sin(2\theta) + \cos(2\theta) = 0 \implies \tan(2\theta) = -\frac{1}{4} = -0.25
\end{equation}
คำนวณมุม $2\theta$ ในจตุรภาคที่สี่ (เพื่อให้ได้ค่าพลังงานต่ำสุด)
\begin{equation}
2\theta = -\arctan(0.25) \approx -0.24498\text{ rad} \implies \theta_{\text{min}} \approx -0.12249\text{ rad} \ (-7.02^\circ)
\end{equation}
คำนวณพลังงานต่ำสุดที่ได้
\begin{equation}
E_{\text{min}} = -2\cos(-0.24498) + 0.5\sin(-0.24498) \approx -2(0.97014) + 0.5(-0.24254) \approx -2.06155\text{ หน่วย}
\end{equation}

\vspace{4pt}
\noindent\textbf{T - Takeaway \& Modern Context (ข้อคิดสรุปและบริบทสมัยใหม่)} \\
ค่าพลังงานที่ได้ $-2.06155$ ตรงกับค่าไอเกนต่ำสุดที่แท้จริงของเมทริกซ์ฮามิลโทเนียนอย่างแม่นยำ ซึ่งพิสูจน์ว่าอันซัทซ์ที่เลือกสามารถครอบคลุมปริภูมิย่อยของสถานะพื้นได้จริง VQE เป็นเทคโนโลยีแกนหลักที่บริษัทยายักษ์ใหญ่กำลังใช้ในการจำลองปฏิกิริยาเคมีเพื่อลดระยะเวลาการสังเคราะห์ยาใหม่
\end{psctbox}

% ตัวอย่างที่ 10.6
\begin{psctbox}{ตัวอย่างที่ 10.6 ---- การตรวจจับความผิดพลาดบิตฟลิปในรหัสสเตบิไลเซอร์ 3 คิวบิต}
\noindent\textbf{P - Problem Statement (โจทย์ปัญหา)} \\
ในรหัสแก้ความผิดพลาดบิตฟลิป 3 คิวบิต สถานะตรรกะนิยามโดย $\ket{0_L} = \ket{000}$ และ $\ket{1_L} = \ket{111}$ ตัวดำเนินการสเตบิไลเซอร์คือ $S_1 = Z_1 Z_2$ และ $S_2 = Z_2 Z_3$ หากคิวบิตที่ 2 เกิดความผิดพลาดบิตฟลิป ($X_2$) จงคำนวณค่าเฉพาะของกลุ่มอาการความผิดพลาด (Error Syndrome) ที่วัดได้ และระบุตำแหน่งของความผิดพลาด

\vspace{4pt}
\noindent\textbf{S - Strategy \& Principle (กลยุทธ์และหลักการ)} \\
พิจารณาสถานะตรรกะทั่วไป $\ket{\psi_L} = \alpha\ket{000} + \beta\ket{111}$ เมื่อถูกรบกวนด้วยเกต $X_2$ คำนวณค่าไอเกนของการวัดตัวดำเนินการสเตบิไลเซอร์ $S_1$ และ $S_2$

\vspace{4pt}
\noindent\textbf{C - Calculation \& Proof (การคำนวณและการพิสูจน์)} \\
สถานะหลังเกิดความผิดพลาดบนคิวบิตที่ 2
\begin{equation}
\ket{\psi_{\text{err}}} = X_2 \ket{\psi_L} = \alpha\ket{010} + \beta\ket{101}
\end{equation}
ทำการวัดตัวดำเนินการสเตบิไลเซอร์ตัวที่หนึ่ง $S_1 = Z_1 Z_2$
\begin{align}
S_1 \ket{010} &= (Z_1\ket{0})(Z_2\ket{1})(I\ket{0}) = (+1)(-1)\ket{010} = -\ket{010} \\
S_1 \ket{101} &= (Z_1\ket{1})(Z_2\ket{0})(I\ket{1}) = (-1)(+1)\ket{101} = -\ket{101}
\end{align}
ดังนั้น $S_1 \ket{\psi_{\text{err}}} = -1 \ket{\psi_{\text{err}}}$ ได้ค่าเฉพาะของ $S_1$ เป็น \textbf{$-1$} \\
ทำการวัดตัวดำเนินการสเตบิไลเซอร์ตัวที่สอง $S_2 = Z_2 Z_3$
\begin{align}
S_2 \ket{010} &= (I\ket{0})(Z_2\ket{1})(Z_3\ket{0}) = (-1)(+1)\ket{010} = -\ket{010} \\
S_2 \ket{101} &= (I\ket{1})(Z_2\ket{0})(Z_3\ket{1}) = (+1)(-1)\ket{101} = -\ket{101}
\end{align}
ดังนั้น $S_2 \ket{\psi_{\text{err}}} = -1 \ket{\psi_{\text{err}}}$ ได้ค่าเฉพาะของ $S_2$ เป็น \textbf{$-1$}

สรุปผลกลุ่มอาการความผิดพลาด $(s_1, s_2) = (-1, -1)$ ซึ่งชี้ชัดว่าความผิดพลาดเกิดขึ้นบนคิวบิตที่ 2 (เนื่องจากคิวบิตที่ 2 อยู่ในทั้ง $S_1$ และ $S_2$) ผู้ควบคุมจึงสามารถแก้ความผิดพลาดได้โดยการกระทำเกต $X_2$ กลับคืนสู่สถานะเดิม

\vspace{4pt}
\noindent\textbf{T - Takeaway \& Modern Context (ข้อคิดสรุปและบริบทสมัยใหม่)} \\
สังเกตว่าค่าเฉพาะที่วัดได้เป็น $-1$ ทั้งสองตัว โดยที่สารสนเทศแอมพลิจูด $\alpha$ และ $\beta$ ไม่ได้ถูกรบกวนเลยแม้แต่น้อย นี่คือความอัศจรรย์ของกรอบสเตบิไลเซอร์ (Stabilizer Formalism) ซึ่งเป็นรากฐานของเซอร์เฟสโค้ดในคอมพิวเตอร์ควอนตัมยุคปัจจุบัน
\end{psctbox}

% ตัวอย่างที่ 10.7
\begin{psctbox}{ตัวอย่างที่ 10.7 ---- ความไวการวัดของเซนเซอร์ควอนตัมที่ขีดจำกัดไฮเซนเบิร์ก}
\noindent\textbf{P - Problem Statement (โจทย์ปัญหา)} \\
ในการวัดเฟส $\phi$ ด้วยอินเทอร์เฟอโรเมทรี หากใช้สถานะพัวพัน NOON ของโฟตอน $N = 100$ อนุภาค $\ket{\text{NOON}} = \frac{1}{\sqrt{2}}(\ket{N, 0} + \ket{0, N})$ จงคำนวณอัตราส่วนการปรับปรุงความไวการวัดเมื่อเปรียบเทียบกับขีดจำกัดมาตรฐานควอนตัม (SQL)

\vspace{4pt}
\noindent\textbf{S - Strategy \& Principle (กลยุทธ์และหลักการ)} \\
ความไม่แน่นอนของเฟสในขีดจำกัดมาตรฐานควอนตัมคือ $\Delta \phi_{\text{SQL}} = 1/\sqrt{N}$ ในขณะที่ในสถานะพัวพัน NOON ความไม่แน่นอนจะบรรลุขีดจำกัดไฮเซนเบิร์ก $\Delta \phi_{\text{Heisenberg}} = 1/N$ คำนวณอัตราส่วนความไว

\vspace{4pt}
\noindent\textbf{C - Calculation \& Proof (การคำนวณและการพิสูจน์)} \\
คำนวณความไม่แน่นอนตามขีดจำกัดมาตรฐานควอนตัม (ไม่มีความพัวพัน)
\begin{equation}
\Delta \phi_{\text{SQL}} = \frac{1}{\sqrt{100}} = \frac{1}{10} = 0.100\text{ rad}
\end{equation}
คำนวณความไม่แน่นอนตามขีดจำกัดไฮเซนเบิร์ก (มีความพัวพันสมบูรณ์)
\begin{equation}
\Delta \phi_{\text{Heisenberg}} = \frac{1}{100} = 0.010\text{ rad}
\end{equation}
อัตราส่วนการปรับปรุงความไวการวัด (Sensitivity Improvement Factor)
\begin{equation}
\text{Improvement} = \frac{\Delta \phi_{\text{SQL}}}{\Delta \phi_{\text{Heisenberg}}} = \frac{1/\sqrt{N}}{1/N} = \sqrt{N} = \sqrt{100} = 10\text{ เท่า}
\end{equation}

\vspace{4pt}
\noindent\textbf{T - Takeaway \& Modern Context (ข้อคิดสรุปและบริบทสมัยใหม่)} \\
ความพัวพันช่วยเพิ่มความไวในการวัดขึ้น 10 เท่าสำหรับ 100 อนุภาค และหากใช้ $N = 10^6$ อนุภาค ความไวจะเพิ่มขึ้นถึง 1,000 เท่า หลักการนี้ถูกนำไปใช้ในเครื่องตรวจวัดคลื่นความโน้มถ่วง LIGO ซึ่งใช้แสงบีบอัดควอนตัม (Quantum Squeezed Light) เพื่อเพิ่มขีดความสามารถในการตรวจจับการชนกันของหลุมดำในห้วงอวกาศลึก
\end{psctbox}

% ==========================================================
% สรุปสาระสำคัญประจำบทที่ 10
% ==========================================================
\section{สรุปสาระสำคัญประจำบทที่ 10}

\begin{itemize}
    \item \textbf{ความได้เปรียบเชิงควอนตัม} ---- เกิดขึ้นจากการผสานพลังของการประมวลผลขนานควอนตัมและการแทรกสอดเชิงควอนตัม เพื่อขยายแอมพลิจูดของคำตอบที่ถูกต้องและหักล้างคำตอบที่ผิดพลาด
    \item \textbf{อัลกอริทึมดอยช์-โยซา} ---- พิสูจน์ความได้เปรียบแบบก้าวกระโดดเชิงกำหนด (Deterministic Exponential Speedup) โดยสามารถแยกแยะฟังก์ชันคงตัวและสมดุลได้ด้วยการเรียกออราเคิลเพียงครั้งเดียวผ่านกลไกการสะท้อนเฟส (Phase Kickback)
    \item \textbf{อัลกอริทึมของโกรเวอร์} ---- ให้การเร่งความเร็วแบบยกกำลังสอง $O(\sqrt{N/M})$ สำหรับการค้นหาในฐานข้อมูลไร้การจัดหมวดหมู่ โดยอาศัยการหมุนเวกเตอร์ในระนาบสองมิติผ่านออราเคิลกลับเฟสและตัวดำเนินการแพร่กระจาย
    \item \textbf{การแปลงฟูเรียร์ควอนตัม (QFT)} ---- ทำงานด้วยความซับซ้อนของเกตเพียง $O(n^2)$ ซึ่งเร็วกว่าขั้นตอนวิธี FFT คลาสสิกแบบเอกซ์โพเนนเชียล เป็นฐานรากของการประเมินเฟสควอนตัม
    \item \textbf{อัลกอริทึมของชอร์} ---- ลดรูปปัญหาการแยกตัวประกอบเฉพาะสู่การหาคาบของฟังก์ชันมอดุโล ทำงานในเวลาพหุนาม $O((\log N)^3)$ ส่งผลให้ต้องมีการเปลี่ยนผ่านสู่มาตรฐานรหัสลับยุคหลังควอนตัม (PQC)
    \item \textbf{ขั้นตอนวิธีในยุค NISQ} ---- VQE และ QAOA ผสานการทำงานระหว่าง QPU และเครื่องปรับปรุงคลาสสิก เพื่อแก้ปัญหาเคมีควอนตัมและการหาค่าเหมาะสมที่สุดโดยไม่ต้องรอคอมพิวเตอร์ควอนตัมที่ทนทานต่อความผิดพร่อง
    \item \textbf{รหัสแก้ความผิดพลาดควอนตัม (QEC)} ---- เซอร์เฟสโค้ดเป็นสถาปัตยกรรมชั้นนำที่ตรวจวัดความผิดพลาดผ่านตัวดำเนินการสเตบิไลเซอร์โดยไม่ทำลายสารสนเทศควอนตัม โดยมีเกณฑ์ทนทานความผิดพลาดประมาณร้อยละ 1
    \item \textbf{ฮาร์ดแวร์และเซนเซอร์ควอนตัม} ---- คิวบิตตัวนำยิ่งยวด ไอออนดักจับ และอะตอมเป็นกลางเป็นเทคโนโลยีหลักในการสร้าง QPU ในขณะที่เซนเซอร์ควอนตัม NV Center เปิดมิติใหม่ในการตรวจวัดระดับนาโนเมตรสำหรับการเกษตรดิจิทัลและสิ่งแวดล้อม
\end{itemize}

% ==========================================================
% แบบฝึกหัดท้ายบทที่ 10
% ==========================================================
\section*{แบบฝึกหัดท้ายบทที่ 10}
\addcontentsline{toc}{section}{แบบฝึกหัดท้ายบทที่ 10}

\subsection*{คำถามเชิงแนวคิด (Conceptual Questions)}

\begin{enumerate}[label=\textbf{\arabic*.}]
    \item จงอธิบายความหมายของคำว่า "ความได้เปรียบเชิงควอนตัม" (Quantum Advantage) และเปรียบเทียบกับคำว่า "อภิมหาอำนาจควอนตัม" (Quantum Supremacy)
    \item เหตุใดการประมวลผลแบบขนานเชิงควอนตัม (Quantum Parallelism) เพียงอย่างเดียวจึงยังไม่เพียงพอที่จะสร้างความได้เปรียบเหนือคอมพิวเตอร์คลาสสิก
    \item อธิบายบทบาทของการแทรกสอดเชิงควอนตัมในการเลือกเฉพาะคำตอบที่ถูกต้องในอัลกอริทึมควอนตัม
    \item กลไกการสะท้อนเฟส (Phase Kickback) ทำงานอย่างไร และเหตุใดจึงจำเป็นต้องเตรียมคิวบิตช่วยในสถานะ $\ket{-}$
    \item ในอัลกอริทึมดอยช์-โยซา เหตุใดผลการวัดคิวบิตแรกเป็น $\ket{0}^{\otimes n}$ จึงยืนยันได้อย่างแน่นอนว่าฟังก์ชันเป็นแบบคงตัว
    \item อธิบายความหมายเชิงเรขาคณิตของการดำเนินการในอัลกอริทึมของโกรเวอร์ในระนาบสองมิติ
    \item ตัวดำเนินการแพร่กระจายของโกรเวอร์ $U_s = 2\ket{s}\bra{s}-I$ ทำหน้าที่ "พลิกกลับรอบค่าเฉลี่ย" (Inversion about the Average) อย่างไร
    \item จะเกิดอะไรขึ้นหากผู้ใช้งานรันอัลกอริทึมของโกรเวอร์ด้วยจำนวนรอบการทำซ้ำมากกว่าจำนวนรอบที่เหมาะสม (Over-rotation)
    \item อธิบายความแตกต่างระหว่างการแปลงฟูเรียร์แบบเร็ว (FFT) และการแปลงฟูเรียร์ควอนตัม (QFT) ในแง่ของความซับซ้อนและการเข้าถึงข้อมูลผลลัพธ์
    \item เหตุใดอัลกอริทึมของชอร์จึงสามารถสั่นคลอนความปลอดภัยของระบบรหัสลับ RSA และ ECC ได้อย่างสิ้นเชิง
    \item อธิบายบทบาทของขั้นตอนวิธีการประเมินเฟสควอนตัม (QPE) ในอัลกอริทึมของชอร์
    \item มาตรฐานวิทยาการรหัสลับยุคหลังควอนตัม (PQC) ใช้หลักการทางคณิตศาสตร์ใดในการป้องกันการโจมตีจากคอมพิวเตอร์ควอนตัม
    \item นิยามของยุค NISQ คืออะไร และเหตุใดอัลกอริทึมแบบผสมผสานคลาสสิก-ควอนตัมจึงเหมาะสมที่สุดสำหรับยุคนี้
    \item อธิบายหลักการทำงานของอัลกอริทึม VQE และบทบาทของทฤษฎีบทแปรผันของเรย์ลีห์-ริทซ์
    \item อันซัทซ์ (Ansatz) ในอัลกอริทึม VQE มีความสำคัญอย่างไร และหากเลือกอันซัทซ์ไม่เหมาะสมจะเกิดผลเสียอย่างไร
    \item อุปสรรคพื้นฐานสามประการในการสร้างรหัสแก้ความผิดพลาดควอนตัม (QEC) เมื่อเปรียบเทียบกับรหัสคลาสสิกคืออะไร
    \item อธิบายหลักการทำงานของเซอร์เฟสโค้ด (Surface Codes) และเหตุใดจึงได้รับความนิยมสูงสุดในการสร้างชิปควอนตัมจริง
    \item เปรียบเทียบข้อดีและข้อจำกัดระหว่างคิวบิตตัวนำยิ่งยวดทรานสมอนและไอออนในกับดักแม่เหล็ก
    \item อธิบายความแตกต่างระหว่างขีดจำกัดมาตรฐานควอนตัม (SQL) และขีดจำกัดของไฮเซนเบิร์กในมาตรวิทยาควอนตัม
    \item จุดบกพร่องไนโตรเจน-ช่องว่างในเพชร (NV Center) สามารถนำมาประยุกต์ใช้ในการเกษตรแม่นยำสูงได้อย่างไร
\end{enumerate}

\subsection*{โจทย์ปัญหาเชิงวิเคราะห์และคำนวณ (Analytical Problems)}

\begin{enumerate}[label=\textbf{\arabic*.}]
    \item ในอัลกอริทึมดอยช์สำหรับ 1 คิวบิต หากออราเคิลแทนฟังก์ชันคงตัว $f(x) = 1$ จงแสดงขั้นตอนการคำนวณเวกเตอร์สถานะและพิสูจน์ว่าผลการวัดคิวบิตแรกจะได้ $\ket{0}$ เสมอ
    \item สำหรับฐานข้อมูลขนาด $N = 64$ รายการ โดยมีคำตอบเป้าหมาย $M = 1$ รายการ จงคำนวณมุม $\theta$ จำนวนรอบการทำซ้ำที่เหมาะสม $R$ และความน่าจะเป็นของความสำเร็จ $P_{\text{success}}$
    \item หากในข้อ 2 ฐานข้อมูลมีคำตอบเป้าหมายเพิ่มขึ้นเป็น $M = 4$ รายการ จงคำนวณจำนวนรอบการทำซ้ำที่เหมาะสม $R$ ใหม่ และเปรียบเทียบกับกรณีคำตอบเดี่ยว
    \item จงแสดงว่าตัวดำเนินการแพร่กระจายของโกรเวอร์ $U_s = 2\ket{s}\bra{s}-I$ เป็นตัวดำเนินการเอกภาพ ($U_s^\dagger U_s = I$) และเป็นตัวดำเนินการแฮร์มิเชียน ($U_s^\dagger = U_s$)
    \item จงเขียนเมทริกซ์การแปลงฟูเรียร์ควอนตัมขนาด $4 \times 4$ สำหรับระบบ 2 คิวบิต ($N=4$) และแสดงว่าเมทริกซ์นี้เป็นเมทริกซ์เอกภาพ
    \item คิวบิต 2 ตัวอยู่ในสถานะ $\ket{\psi} = \frac{1}{2}(\ket{00} + \ket{01} + \ket{10} + \ket{11})$ จงคำนวณผลการแปลง $\text{QFT}\ket{\psi}$
    \item ในการแยกตัวประกอบ $N = 21$ โดยใช้อัลกอริทึมของชอร์ หากสุ่มเลือก $a = 2$ จงคำนวณหาคาบ $r$ ของฟังก์ชัน $2^x \bmod 21$ และแสดงการคำนวณตัวประกอบเฉพาะ $p$ และ $q$
    \item ในการประเมินเฟสควอนตัม (QPE) หากต้องการประเมินเฟส $\theta$ ให้มีความแม่นยำ $n = 4$ บิต โดยมีความน่าจะเป็นของความสำเร็จไม่น้อยกว่า $1 - \epsilon = 0.90$ จงคำนวณจำนวนคิวบิตนับขั้นต่ำ $t$ ที่ต้องใช้
    \item กำหนดฮามิลโทเนียนอย่างง่าย $\hat{H} = 2 Z_1 Z_2 + X_1 + X_2$ หากใช้อันซัทซ์ $\ket{\psi(\theta)} = \cos(\theta)\ket{01} + \sin(\theta)\ket{10}$ จงหาฟังก์ชันพลังงาน $E(\theta)$ และคำนวณพลังงานต่ำสุด
    \item ในรหัสแก้ความผิดพลาดเฟสฟลิป 3 คิวบิต สถานะตรรกะคือ $\ket{0_L} = \ket{+++}$ และ $\ket{1_L} = \ket{---}$ ตัวดำเนินการสเตบิไลเซอร์คือ $S_1 = X_1 X_2$ และ $S_2 = X_2 X_3$ หากคิวบิตที่ 3 เกิดความผิดพลาดเฟสฟลิป ($Z_3$) จงคำนวณกลุ่มอาการความผิดพลาดที่วัดได้
    \item จงคำนวณจำนวนคิวบิตกายภาพที่ต้องใช้ในการสร้างเซอร์เฟสโค้ดที่มีระยะทาง $d = 7$ (ซึ่งสามารถแก้ความผิดพลาดได้ $t = (d-1)/2 = 3$ ตัว) หากใช้สถาปัตยกรรมมาตรฐานแบบหมุน ($d^2$ คิวบิตข้อมูล และ $d^2-1$ คิวบิตช่วย)
    \item ในการค้นหาของโกรเวอร์ หากฐานข้อมูลมีขนาด $N = 10^8$ รายการ จงเปรียบเทียบจำนวนการเรียกใช้ออราเคิลระหว่างวิธีคลาสสิกเฉลี่ย ($N/2$) และอัลกอริทึมโกรเวอร์ และคำนวณอัตราเร่งความเร็ว
    \item สำหรับสถานะ GHZ 3 คิวบิต $\ket{\text{GHZ}} = \frac{1}{\sqrt{2}}(\ket{000}+\ket{111})$ จงคำนวณผลการกระทำของการแปลงฟูเรียร์ควอนตัม $\text{QFT}^{\otimes 3} \ket{\text{GHZ}}$
    \item หากเกตในวงจรควอนตัมมีความแม่นยำ $F_{\text{gate}} = 99.8\%$ จงคำนวณความแม่นยำรวมของวงจร (Circuit Fidelity) ที่ประกอบด้วยเกตสองคิวบิตจำนวน 500 เกต โดยสมมติว่าความผิดพลาดเป็นอิสระต่อกัน
    \item เซนเซอร์ควอนตัม NV Center มีค่าอัตราส่วนไจโรแมกเนติกของสปินอิเล็กตรอน $\gamma_e / 2\pi \approx 28\text{ GHz/T}$ หากสนามแม่เหล็กภายนอกเปลี่ยนแปลงไป $\Delta B = 10\ \mu\text{T}$ จงคำนวณการเลื่อนของความถี่การสั่นพ้องซีมาน $\Delta f$
    \item ในการวัดสัญญาณด้วยอินเทอร์เฟอโรเมทรีของอะตอมเย็น หากใช้กลุ่มหมอกอะตอมจำนวน $N = 10^6$ อะตอม จงคำนวณขีดจำกัดความไม่แน่นอนของเฟส $\Delta \phi$ ภายใต้ขีดจำกัดมาตรฐานควอนตัม (SQL) และขีดจำกัดไฮเซนเบิร์ก
    \item พิจารณาวงจรอัลกอริทึม QAOA ขนาด 1 ชั้น ($p=1$) สำหรับกราฟเส้นตรง 2 จุดเชื่อมต่อกันด้วย 1 เส้นเชื่อม จงเขียนสมการตัวดำเนินการวิวัฒน์ $e^{-i\beta \hat{H}_M} e^{-i\gamma \hat{H}_C}$ ในรูปเมทริกซ์
    \item จงแสดงว่าเกตสลับคิวบิต SWAP สามารถสร้างขึ้นจากเกต CNOT จำนวน 3 เกตที่สลับทิศทางควบคุมกันได้อย่างไร
    \item ในการทดสอบความได้เปรียบเชิงควอนตัมด้วยการสุ่มตัวอย่างวงจรควอนตัม (Random Circuit Sampling) จงอธิบายว่าเหตุใดการจำลองการแจกแจงแบบพอร์เตอร์-โธมัส (Porter-Thomas Distribution) จึงเป็นเรื่องยากอย่างยิ่งยวดสำหรับซูเปอร์คอมพิวเตอร์คลาสสิก
    \item \textbf{โจทย์ประยุกต์การคำนวณเชิงตัวเลขด้วย Qiskit (Computational Task)} ---- จงเขียนโค้ดภาษา Python โดยใช้ไลบรารี Qiskit เพื่อสร้างวงจรการแปลงฟูเรียร์ควอนตัมผกผัน ($\text{QFT}^\dagger$) ขนาด 3 คิวบิต จากนั้นนำไปเชื่อมต่อเข้ากับวงจรการประเมินเฟสควอนตัม (QPE) สำหรับตัวดำเนินการเอกภาพ $U = \begin{pmatrix} 1 & 0 \\ 0 & e^{2\pi i (3/8)} \end{pmatrix}$ พร้อมทั้งจำลองผลการวัดและตรวจสอบว่าบิตเอาต์พุตที่ได้ตรงกับค่าฐานสอง $0.011_2$ ($3/8$) อย่างแม่นยำ
\end{enumerate}
